% !TeX root = ../main.tex

\documentclass[../main.tex]{subfiles}
\begin{document}

% ============================================================
\section{Mô hình toán học đối tượng điều khiển}
% ============================================================

Dựa trên phân tích đặc tính động học từ dữ liệu thực nghiệm, hai mô hình toán học được xây dựng để đại diện cho hai đối tượng điều khiển trong hệ điều khiển hai vòng.

Đối với vòng trong, đối tượng van điều khiển lưu lượng có đặc tính đáp ứng nhanh với thời gian trễ nhỏ. Dựa trên hình dạng đáp ứng bậc thang, mô hình QTB1T (First-Order Plus Dead Time) được xác định là phù hợp:
\begin{equation}
    G_v(s) = \frac{K_v}{T_v s + 1}\, e^{-L_v s}
    \label{eq:qtb1t}
\end{equation}

trong đó $K_v$ là hệ số khuếch đại tĩnh, $T_v$ là hằng số thời gian quán tính, và $L_v$ là thời gian trễ vận chuyển. Mô hình này có ba thông số cần xác định thông qua thuật toán nhận dạng tối ưu hóa.

Đối với vòng ngoài, quá trình sinh hơi trong lò hơi có đặc tính đáp ứng chậm với đường cong dạng chữ S rõ rệt. Mô hình QTB2T (Second-Order Plus Dead Time) được lựa chọn để phản ánh đặc tính hai giai đoạn của quá trình truyền nhiệt:
\begin{equation}
    G_p(s) = \frac{K_p}{(T_{1p} s + 1)(T_{2p} s + 1)}\, e^{-L_p s}
    \label{eq:qtb2t}
\end{equation}

trong đó $K_p$ là hệ số khuếch đại, $T_{1p}$ và $T_{2p}$ là hai hằng số thời gian quán tính, và $L_p$ là thời gian trễ. Hai hằng số thời gian phản ánh hai giai đoạn trong quá trình truyền nhiệt: giai đoạn truyền nhiệt từ ngọn lửa đến bề mặt ống (bức xạ), và giai đoạn truyền nhiệt từ bề mặt ống đến nước trong ống (đối lưu).

\begin{figure}[H]
    \centering
    \begin{tikzpicture}[
        node distance=1.5cm and 1.2cm,
        inner block/.style={draw, rectangle, minimum width=4.5cm, minimum height=1.2cm, align=center, thick},
        outer block/.style={draw, rectangle, minimum width=4.5cm, minimum height=1.2cm, align=center, thick},
        arrow/.style={-Latex, thick},
        label style/.style={font=\bfseries\small, anchor=east},
        io/.style={font=\small}
    ]
        %% HÀNG 1: VÒNG TRONG (VAN)
        \node[label style] (lblVan) at (0, 2) {Vòng trong (Van):};
        \node[io, right=0.5cm of lblVan] (uv) {$u_v$};
        \node[inner block, right=of uv] (Gv) {$G_v(s)=\dfrac{K_v}{T_v s+1}\,e^{-L_v s}$};
        \node[io, right=of Gv] (Fv) {$F(s)$};

        %% HÀNG 2: VÒNG NGOÀI (LÒ)
        \node[label style] (lblLo) at (0, 0) {Vòng ngoài (Lò):};
        \node[io] (Fp) at (uv |- lblLo) {$F(s)$};
        \node[inner block, right=-2cm of Fp] (Gp) at (Gv |- lblLo) {$G_p(s)=\dfrac{K_p}{(T_{1p}s+1)(T_{2p}s+1)}\,e^{-L_p s}$};
        \node[io, right=of Gp] (Pp) {$P(s)$};

        %% ĐƯỜNG MŨI TÊN
        \draw[arrow] (uv) -- (Gv);
        \draw[arrow] (Gv) -- (Fv);
        \draw[arrow] (Fp) -- (Gp);
        \draw[arrow] (Gp) -- (Pp); 
        
    \end{tikzpicture}
    \caption{Sơ đồ khối hàm truyền của các đối tượng điều khiển}
    \label{fig:so_do_ham_truyen}
\end{figure}

\begin{figure}[H]
    \centering
\begin{tikzpicture}[
    >=Stealth,
    scale=0.75,
    transform shape
]

% Grid
\draw[gray!30, very thin] (0,0) grid (10,5);

% Axes
\draw[thick, ->] (0,0) -- (10,0) node[below] {Thời gian $t$ (s)};
\draw[thick, ->] (0,0) -- (0,5.5) node[left] {$y(t)$};

% Labels
\node[below left] at (0,0) {$0$};

% Dead time annotation
\draw[dashed, red!60!black, thick] (1.5,0) -- (1.5,4.5);
\draw[<->, red!60!black] (0,-0.3) -- (1.5,-0.3);
\node[red!60!black, below] at (0.75,-0.3) {$L$};

% QTB1T curve
\draw[blue!70!black, very thick, domain=1.5:9.5, smooth] plot ({\x}, {4*(1-exp(-(\x-1.5)/2))});

% Mark K
\draw[dashed, gray!60!black] (9,4) -- (9,0);
\node[right, gray!60!black] at (9.1,4) {$K$};

% Mark T (63.2% of K)
\draw[dashed, green!60!black] (3.5,1.26) -- (3.5,0);
\node[below, green!60!black] at (3.5,0) {$L+T$};

% Time constant line
\draw[<->, green!60!black] (1.5,5.2) -- (3.5,5.2);
\node[above, green!60!black, midway] at (2.5,5.2) {$T$};

% Rise to 63.2%
\fill[green!60!black] (3.5,1.26) circle (2pt);
\node[right, green!60!black] at (3.6,1.26) {\small $0.632K$};

% Step input
\draw[gray!50!black, thick] (0,0) -- (0,4);
\draw[gray!50!black, very thick] (0,4) -- (1.5,4);
\draw[gray!50!black, thick] (1.5,4) -- (1.5,0);
\node[left, gray!50!black] at (0,2) {$u(t)$};

% Legend box -- placed inside the plot area (upper right corner)
\node[draw, fill=yellow!15, rounded corners, anchor=north east, font=\small] at (15.9,4.9) {
    \begin{tabular}{p{3.2cm}}
    \textbf{Chú thích:}\\
    $K$: Hệ số khuếch đại tĩnh\\
    $T$: Hằng số thời gian\\
    $L$: Thời gian trễ\\[4pt]
    $y(t) = K(1-e^{-(t-L)/T})$
    \end{tabular}
};

% Annotation for rise time -- placed inside the plot area (lower right)
\node[anchor=north east, draw, fill=white, font=\small] at (9.9, 3) {
    \begin{tabular}{p{3.2cm}}
    \textbf{Đặc tính QTB1T:}\\
    Đáp ứng bắt đầu tại $t = L$
    \end{tabular}
};

\end{tikzpicture}
    \caption{Đáp ứng bậc thang đơn vị của mô hình QTB1T với các thông số động học}
    \label{fig:qtb_response}
\end{figure}

Hình~\ref{fig:qtb_response} minh họa đáp ứng bậc thang đơn vị với các thông số động học được chú thích:
\begin{itemize}
    \item \textbf{$L$ (Dead time)}: Thời gian trễ ban đầu trước khi đầu ra bắt đầu thay đổi
    \item \textbf{$T$ (Time constant)}: Hằng số thời gian -- thời gian để đầu ra đạt 63.2\% giá trị xác lập
    \item \textbf{$K$ (Static gain)}: Hệ số khuếch đại tĩnh -- giá trị xác lập cuối cùng
\end{itemize}

\begin{figure}[H]
    \centering
\begin{tikzpicture}[
    >=Stealth,
    scale=0.75,
    transform shape
]

% Grid
\draw[gray!30, very thin] (0,0) grid (10,4.5);

% Axes
\draw[thick, ->] (0,0) -- (10,0) node[below] {Thời gian $t$ (s)};
\draw[thick, ->] (0,0) -- (0,5) node[left] {$y(t)$};

% Dead time lines
\draw[dashed, gray] (1,0) -- (1,4.5);

% QTB1T: y = K*(1-exp(-(t-L)/T)) with L=1, T=2.5, K=4
\draw[blue, very thick, domain=1:9.5, smooth] plot ({\x}, {4*(1-exp(-(\x-1)/2.5))});

% QTB2T: y = K*(1 - (T1*exp(-(t-L)/T1) - T2*exp(-(t-L)/T2))/(T1-T2)) with L=1, T1=1.5, T2=0.5, K=4
\draw[red!70!black, very thick, domain=1:9.5, smooth] plot ({\x}, {4*(1-(1.5*exp(-(\x-1)/1.5)-0.5*exp(-(\x-1)/0.5))/(1.5-0.5))});

% Mark inflection point
\fill[red!70!black] (2.5,1.5) circle (2.5pt);
\node[red!70!black, above right] at (2.5,1.5) {\small Điểm uốn};

% Legend -- placed inside plot (upper right)
\node[anchor=north east, draw, fill=white, font=\small] at (13.9, 5) {
    \begin{tabular}{p{3.95cm}}
    \textcolor{blue}{\textbf{---}} QTB1T\\
    \quad $\dfrac{K}{Ts+1}e^{-Ls}$ \\[8pt]
    \textcolor{red!70!black}{\textbf{---}} QTB2T\\
    \quad $\dfrac{K}{(T_1s+1)(T_2s+1)}e^{-Ls}$
    \end{tabular}
};

% Characteristics box -- placed inside plot (middle right)
\node[anchor=north east, draw, fill=blue!5, font=\small] at (18.9, 5.0) {
    \begin{tabular}{p{4cm}}
    \textbf{Đặc điểm QTB2T:}\\
    Đường cong hình chữ S\\
    Có điểm uốn rõ rệt\\
    Quy trình hai giai đoạn\\
    Giai đoạn 1: Bức xạ\\
    Giai đoạn 2: Đối lưu
    \end{tabular}
};

% Step input
\draw[gray!50!black, thick] (0,0) -- (0,4);
\draw[gray!50!black, very thick] (0,4) -- (1,4);
\draw[gray!50!black, thick] (1,4) -- (1,0);
\node[left, gray!50!black] at (0,2) {$u(t)$};

\end{tikzpicture}
    \caption{So sánh đáp ứng bậc thang của mô hình QTB1T và QTB2T}
    \label{fig:qtb1t_vs_qtb2t}
\end{figure}

Hình~\ref{fig:qtb1t_vs_qtb2t} so sánh trực tiếp đáp ứng của hai mô hình:
\begin{itemize}
    \item \textbf{QTB1T}: Đáp ứng tăng theo hàm mũ đơn sau thời gian trễ $L$
    \item \textbf{QTB2T}: Đáp ứng có điểm uốn rõ rệt, phản ánh quá trình hai giai đoạn
\end{itemize}

Việc lựa chọn mô hình QTB1T cho van và QTB2T cho lò hơi dựa trên cơ sở lý thuyết và thực nghiệm. Mô hình QTB1T là xấp xỉ đơn giản nhưng đủ chính xác đối với các đối tượng có đặc tính bậc thấp, trong khi mô hình QTB2T phản ánh tốt hơn đặc tính hai giai đoạn của quá trình nhiệt trong lò hơi. Các thông số cụ thể của mô hình sẽ được xác định bằng thuật toán nhận dạng tối ưu hóa trong phần tiếp theo.

% ============================================================
\section{Xử lý số liệu thực nghiệm}
% ============================================================

Để có cơ sở dữ liệu tin cậy phục vụ cho công tác nhận dạng mô hình toán học và tổng hợp thuật toán điều khiển, nhóm đã sử dụng số liệu thực nghiệm mô tả đáp ứng quá độ của cả hai đối tượng: vòng trong (đối tượng van điều khiển lưu lượng) và vòng ngoài (đối tượng lò hơi sinh áp suất).

Phương pháp thu thập của số liệu được thực hiện bằng cách áp dụng một tín hiệu kích thích dạng bậc thang (Step Input) vào đầu vào của hệ thống khi đang ở trạng thái làm việc ổn định, sau đó sử dụng hệ thống thu thập dữ liệu (DAQ) để ghi nhận diễn biến của tín hiệu đầu ra theo thời gian. Các tập dữ liệu thu được lưu trữ dưới định dạng tệp CSV, cụ thể là tệp {\footnotesize\texttt{van-dap-ung-xung-1,6t.csv}} đối với đáp ứng của lưu lượng dầu, và tệp đáp ứng áp suất hơi tương ứng. Hình~\ref{fig:preprocessing} minh họa quy trình tiền xử lý dữ liệu thực nghiệm.

\begin{figure}[H]
    \centering
    \begin{tikzpicture}[
        >=Stealth,
        scale=0.7,
        transform shape
    ]

    % Flowchart boxes
    \node[draw, rectangle, fill=blue!10, minimum width=4cm, align=center] (raw) at (0,0) {Dữ liệu thô\\({\footnotesize\texttt{van-dap-ung-xung-1,6t.csv}})};

    \node[draw, diamond, fill=yellow!20, minimum size=1.5cm, align=center] (filter) at (0,-2.5) {Lọc nhiễu};

    \node[draw, rectangle, fill=yellow!20, minimum width=4cm, align=center] (filtered) at (0,-5) {Dữ liệu đã lọc\\Loại bỏ nhiễu tần số cao};

    \node[draw, diamond, fill=orange!20, minimum size=1.5cm, align=center] (outlier) at (0,-8) {Phát hiện\\outlier};

    \node[draw, rectangle, fill=orange!20, minimum width=4cm, align=center] (clean) at (0,-11) {Dữ liệu sạch\\Chuẩn bị nhận dạng};

    % Arrows
    \draw[->, thick] (raw) -- (filter);
    \draw[->, thick] (filter) -- (filtered);
    \draw[->, thick] (filtered) -- (outlier);
    \draw[->, thick] (outlier) -- (clean);

    % Labels
    \node[right=0.5cm of filter, align=center] {Moving Average\\(cửa sổ 5 mẫu)};
    \node[right=0.5cm of outlier, align=center] {IQR Method};

    % Process description
    \node[anchor=west, draw, fill=white] at (4,-5) {
        \begin{tabular}{p{5.5cm}}
        \textbf{Thông số tiền xử lý:}\\
        • Cửa sổ lọc: 5 mẫu\\
        • Tần số lấy mẫu: $f_s = 10$ Hz\\
        • Ngưỡng outlier: $1.5 \times$ IQR
        \end{tabular}
    };

    \end{tikzpicture}
    \caption{Sơ đồ quy trình tiền xử lý dữ liệu thực nghiệm}
    \label{fig:preprocessing}
\end{figure}

Bảng \ref{tab:so_lieu_van} và Bảng \ref{tab:so_lieu_ap_suat} trình bày một phần dữ liệu đã được trích xuất từ các tệp thực nghiệm, minh họa rõ nét xu hướng tăng dần của đáp ứng đầu ra sau khi nhận được tín hiệu kích thích.

\begin{table}[htbp]
\centering
\caption{Trích xuất số liệu thực nghiệm đáp ứng van điều khiển lưu lượng}
\label{tab:so_lieu_van}
\scalebox{0.75}{
\begin{tabular}{|c|c|c|c|}
\hline
\textbf{Thời gian (s)} & \textbf{Đáp ứng lưu lượng (t/h)} & \textbf{Thời gian (s)} & \textbf{Đáp ứng lưu lượng (t/h)}\\ \hline
0.00 & 0.000 & 2.40 & 0.379 \\ \hline
0.80 & 0.001 & 5.52 & 0.856 \\ \hline
1.17 & 0.089 & 10.49 & 1.203 \\ \hline
1.54 & 0.175 & 29.77 & 1.399 \\ \hline
\end{tabular}
}
\end{table}

\begin{table}[htbp]
\centering
\caption{Trích xuất số liệu thực nghiệm đáp ứng áp suất hơi}
\label{tab:so_lieu_ap_suat}
\scalebox{0.78}{
\begin{tabular}{|c|c|c|c|}
\hline
\textbf{Thời gian (s)} & \textbf{Đáp ứng áp suất (bar)} & \textbf{Thời gian (s)} & \textbf{Đáp ứng áp suất (bar)}\\ \hline
0.00 & $-$0.002 & 15.02 & 0.660 \\ \hline
1.69 & $-$0.002 & 20.14 & 0.929 \\ \hline
3.76 & 0.028 & 25.31 & 1.087 \\ \hline
7.17 & 0.140 & 48.04 & 1.197 \\ \hline
\end{tabular}
}
\end{table}

Đặc điểm của các tập số liệu này thể hiện tính chất quá độ điển hình của các đối tượng công nghiệp: có thời gian trễ ban đầu và đường cong đáp ứng hình chữ S tiệm cận dần về giá trị xác lập.

% \begin{figure}[H]
%     \centering
%     \begin{tikzpicture}[
%         >=Stealth,
%         scale=0.85,
%         transform shape
%     ]

%     % Left Box: Raw Data
%     \node[draw, rectangle, fill=red!5, minimum width=4.5cm, minimum height=2.5cm, align=center, rounded corners] (raw) at (-4.5, 0) {
%         \begin{tabular}{p{6cm}}
%         \textbf{Dữ liệu thô:}\\
%         • Nhiễu ngẫu nhiên cao tần\\
%         • Chứa các giá trị ngoại lai\\
%         • Gây sai số khi nhận dạng
%         \end{tabular}
%     };

%     % Right Box: Clean Data
%     \node[draw, rectangle, fill=green!5, minimum width=4.5cm, minimum height=2.5cm, align=center, rounded corners] (clean) at (4.5, 0) {
%         \begin{tabular}{p{6cm}}
%         \textbf{Dữ liệu đã lọc:}\\
%         • Đã loại bỏ nhiễu cao tần\\
%         • Loại bỏ các outlier (IQR)\\
%         • Giữ nguyên đặc tính động học
%         \end{tabular}
%     };

%     % Transformation Arrow
%     \draw[->, very thick, blue!70!black] (raw) -- (clean) 
%         node[midway, above, font=\bfseries\small, text=blue!70!black] {Tiền xử lý};

%     \end{tikzpicture}
%     \caption{So sánh dữ liệu thô và dữ liệu sau tiền xử lý}
%     \label{fig:data_comparison}
% \end{figure}

% Hình~\ref{fig:data_comparison} thể hiện dữ liệu trước và sau khi tiền xử lý. Phần lớn nhiễu ngẫu nhiên đã được loại bỏ trong khi hình dạng đặc tính động học chính được bảo toàn.

% \begin{figure}[H]
%     \centering
%     \includegraphics[width=0.7\textwidth]{Hinhve/dac-tinh-van-dap-ung-xung-1,6t.png}
%     \caption{Đáp ứng bậc thang thực nghiệm của đối tượng van điều khiển lưu lượng dầu}
%     \label{fig:van_dap_ung_xung_2}
% \end{figure}

% Hình~\ref{fig:van_dap_ung_xung_2} thể hiện đáp ứng bậc thang gốc (raw data) của đối tượng van, cho thấy dữ liệu sau khi tiền xử lý đã loại bỏ phần lớn nhiễu ngẫu nhiên và giữ được hình dạng đặc tính động học chuẩn xác.

% Nguồn số liệu mang tính thực tiễn cao này sẽ là dữ liệu đầu vào trực tiếp cho thuật toán nhận dạng phi tuyến, giúp xác định các thông số động học chính xác cho mô hình QTB1T và QTB2T.

% ============================================================
\section{Nhận dạng đối tượng điều khiển}
% ============================================================

Trong điều khiển tự động hóa hiện đại, quá trình tổng hợp và chỉnh định bộ điều khiển phụ thuộc hoàn toàn vào mức độ chính xác của mô hình toán học đại diện cho đối tượng thực tế. Phương pháp nhận dạng đối tượng được lựa chọn dựa trên việc phân tích đáp ứng bậc thang của hệ thống trong miền thời gian.

Về mặt toán học, quá trình nhận dạng được định hình thành bài toán tối ưu hóa phi tuyến: tìm kiếm tập tham số $A$ sao cho hàm mục tiêu sai số toàn phương $J$ đạt giá trị cực tiểu:

\begin{figure}[!htbp]
    \centering
    \begin{tikzpicture}[
        >=Stealth,
        scale=0.75,
        transform shape,
        flowchart/.style={draw, rectangle, minimum width=3cm, minimum height=1cm, align=center}
    ]

    % Start
    \node[flowchart, fill=green!20] (init) at (0,0) {Khởi tạo tham số\\$A_0 = [K_0, T_0, L_0]$};

    % Calculate model
    \node[flowchart, fill=blue!10, below=0.6cm of init] (calc) {Tính $y_{model}(t_i, A)$};

    % Calculate error
    \node[flowchart, fill=blue!10, below=0.6cm of calc] (error) {Tính sai số $e_i$};

    % Calculate objective
    \node[flowchart, fill=yellow!20, below=0.6cm of error] (obj) {Tính hàm mục tiêu $J(A)$};

    % Check convergence
    \node[draw, diamond, fill=orange!20, below=0.6cm of obj, minimum size=1.5cm, aspect=2] (check) {$J < \epsilon$?};

    % Update parameters
    \node[flowchart, fill=blue!10, right=1.5cm of obj] (update) {Cập nhật\\$A \leftarrow A + \Delta A$};

    % Result
    \node[flowchart, fill=green!20, below=0.8cm of check, minimum width=3.5cm] (result) {Tham số tối ưu\\$A^* = [K^*, T^*, L^*]$};

    % Arrows
    \draw[->, thick] (init) -- (calc);
    \draw[->, thick] (calc) -- (error);
    \draw[->, thick] (error) -- (obj);
    \draw[->, thick] (obj) -- (check);
    \draw[->, thick] (check.east) -| node[pos=0.25, above] {Không} (update.south);
    \draw[->, thick] (check.south) -- node[right] {Có} (result.north);
    \draw[->, thick] (update.north) |- (calc.east);

    % Legend
    \node[anchor=north west, draw, fill=white] at (8, -2) {
        \begin{tabular}{p{5cm}}
        \textbf{Thuật toán: Grid Search}\\[4pt]
        • Không cần tính đạo hàm\\
        • Quét lưới 3--4 biến\\
        • Chạy nhiều lần với các\\ \quad điểm khởi tạo khác nhau
        \end{tabular}
    };

    \end{tikzpicture}
    \caption{Sơ đồ thuật toán nhận dạng tối ưu hóa}
    \label{fig:identification_algorithm}
\end{figure}

Một bước tiền xử lý quan trọng trước khi nhận dạng là chuẩn hóa dữ liệu. Do tín hiệu bậc thang kích thích có biên độ khác nhau (biên độ 1.6 đối với van và 1.0 đối với áp suất hơi), dữ liệu được chuẩn hóa bằng cách chia cho biên độ bậc thang tương ứng, đảm bảo các tham số nhận dạng ($K$, $T$, $L$) phản ánh đặc tính động học thuần túy.

Để đánh giá chất lượng nhận dạng, cần sử dụng thêm hệ số xác định $R^2$ (càng gần 1 càng tốt) và sai số bình phương trung bình gốc RMSE (càng nhỏ càng tốt):
\begin{equation}
    R^2 = 1 - \frac{\sum_{i=1}^{N}(y_i - \hat{y}_i)^2}{\sum_{i=1}^{N}(y_i - \bar{y})^2}
\end{equation}
\begin{equation}
    \text{RMSE} = \sqrt{\frac{1}{N}\sum_{i=1}^{N}(y_i - \hat{y}_i)^2}
\end{equation}

Lựa chọn thuật toán tối ưu hóa bình phương tối thiểu kết hợp quét lưới (grid search) thay vì các phương pháp nhận dạng truyền thống vì: quét lưới không cần tính đạo hàm, đơn giản trong triển khai, và hiệu quả đối với bài toán 3--4 biến. Thuật toán được chạy nhiều lần với các điểm khởi tạo khác nhau, kết quả cuối cùng được chọn từ lần chạy cho giá trị $J$ nhỏ nhất.

Với cấu trúc điều khiển hai vòng đã đề xuất, hai khâu động học cần được nhận dạng độc lập: đối tượng vòng trong (cụm van điều khiển lưu lượng dầu) và đối tượng vòng ngoài (quá trình sinh áp suất hơi của lò).

\subsection{Nhận dạng đối tượng van (Vòng trong)}

Đối tượng van điều khiển có đặc tính động học tương đối nhanh nhưng vẫn tồn tại độ trễ do cơ khí và ma sát đường ống dẫn. Do đó, mô hình xấp xỉ quán tính bậc nhất có trễ (QTB1T) được lựa chọn. Hàm truyền có dạng như trong Phương trình~\ref{eq:qtb1t}.

Đáp ứng bậc thang đơn vị trong miền thời gian:
\begin{equation}
    y_{\mathrm{model}}(t) =
    \begin{cases}
        0, & t < L_v \\[4pt]
        K_v\!\left(1 - e^{-\frac{t - L_v}{T_v}}\right), & t \geq L_v
    \end{cases}
    \label{eq:qtb1t_step}
\end{equation}

Tiến hành đưa dữ liệu thực nghiệm từ tệp {\footnotesize\texttt{van-dap-ung-xung-1,6t.csv}} (biên độ tín hiệu bậc thang đầu vào $1.6$) vào thuật toán, hệ thống hội tụ và xuất ra bộ thông số tối ưu:

\begin{itemize}
    \item Hệ số khuếch đại tĩnh: $K_v = 0{,}8836$
    \item Hằng số thời gian quán tính: $T_v = 4{,}8694$ s
    \item Thời gian trễ vận chuyển: $L_v = 0{,}8213$ s
\end{itemize}

Giá trị $K_v = 0{,}8836$ cho thấy hệ thống có hệ số khuếch đại nhỏ hơn 1, phản ánh thực tế độ mở van tối đa không tương ứng với lưu lượng danh nghĩa thiết kế do đặc tính thủy lực và sụt áp trên đường ống. Giá trị $T_v = 4{,}8694$ s phản ánh tốc độ đáp ứng nhanh của van. Giá trị $L_v = 0{,}8213$ s phản ánh tổng thời gian trễ bao gồm trễ vận chuyển trên đường ống dầu và trễ do đáp ứng cơ khí của van.

Hệ số $R^2 = 0{,}9987$ cho thấy mô hình QTB1T giải thích được 99.87\% biến thiên của dữ liệu thực nghiệm, xác nhận cấu trúc QTB1T hoàn toàn phù hợp với đặc tính động học của van.

\begin{table}[H]
\centering
\caption{Kết quả nhận dạng đối tượng van điều khiển lưu lượng}
\label{tab:ket_qua_nhan_dang_van}
\scalebox{0.8}{
\begin{tabular}{|l|c|}
\hline
\textbf{Thông số} & \textbf{Giá trị}\\ \hline
Hệ số khuếch đại $K_v$ & $0{,}8836$\\ \hline
Hằng số thời gian $T_v$ (s) & $4{,}8694$\\ \hline
Thời gian trễ $L_v$ (s) & $0{,}8213$\\ \hline
Hệ số xác định $R^2$ & $0{,}9987$\\ \hline
Sai số RMSE & $0{,}0142$\\ \hline
\end{tabular}
}
\end{table}

\begin{figure}[H]
    \centering
    \includegraphics[width=0.8\textwidth]{Hinhve/nhan_dang_van.png}
    \caption{Đồ thị so sánh kết quả nhận dạng đối tượng van điều khiển theo mô hình QTB1T}
    \label{fig:nhan_dang_van}
\end{figure}

Biểu đồ tại Hình~\ref{fig:nhan_dang_van} cho thấy đường cong mô phỏng bám sát hoàn toàn với tập dữ liệu thực nghiệm, minh chứng cho tính đúng đắn và độ chính xác của cấu trúc QTB1T đối với phân hệ van lưu lượng.

% \begin{figure}[!htbp]
%     \centering
% \begin{tikzpicture}[>=Stealth, scale=0.75, transform shape]

% \begin{axis}[
%     width=0.4\textwidth,
%     height=5cm,
%     at={(0,0)},
%     anchor=west,
%     xlabel={Thời gian (s)},
%     ylabel={Phần dư},
%     title={Đồ thị phần dư theo thời gian},
%     grid=major,
%     legend pos=north east
% ]
% \addplot[blue, only marks, mark=*, mark size=1.5] coordinates {
%     (0, 0.002) (1, -0.003) (2, 0.001) (3, -0.002) (4, 0.003)
%     (5, -0.001) (6, 0.002) (7, -0.003) (8, 0.001) (9, -0.002)
%     (10, 0.002) (11, -0.001) (12, 0.003) (13, -0.002) (14, 0.001)
% };
% \addlegendentry{Phần dư}
% \addplot[red!50!black, dashed, domain=0:14] {0.01};
% \addplot[red!50!black, dashed, domain=0:14] {-0.01};
% \addlegendentry{$\pm3\sigma$}
% \addplot[gray!70!black, dashed, domain=0:14] {0};
% \addlegendentry{0}
% \end{axis}

% \begin{axis}[
%     width=0.4\textwidth,
%     height=5cm,
%     at={(0.5\textwidth,0)},
%     anchor=west,
%     xlabel={Phần dư chuẩn hóa},
%     ylabel={Tần suất},
%     title={Histogram phần dư},
%     grid=major,
%     legend pos=north east,
%     ybar interval=1
% ]
% \addplot[blue!70!black, fill=blue!30] coordinates {
%     (-0.015, 1) (-0.01, 2) (-0.005, 4) (0, 6) (0.005, 4) (0.01, 2) (0.015, 1)
% };
% \addlegendentry{Tần suất}
% \end{axis}

% \end{tikzpicture}
%     \caption{Kiểm chứng mô hình van: phân tích phần dư (Residual Analysis)}
%     \label{fig:validation_van}
% \end{figure}

% Hình~\ref{fig:validation_van} thể hiện kết quả kiểm chứng mô hình van bằng phương pháp phân tích phần dư:
% \begin{itemize}
%     \item Phần dư dao động ngẫu nhiên quanh giá trị 0, không có xu hướng hệ thống
%     \item Phân bố phần dư tuân theo phân bố chuẩn (xấp xỉ Gaussian)
%     \item Không có autocorrelation đáng kể (xác nhận independence)
% \end{itemize}

\subsection{Nhận dạng đối tượng áp suất hơi (Vòng ngoài)}

Quá trình sinh hơi trong lò mang bản chất tích lũy năng lượng và trao đổi nhiệt phức tạp qua nhiều bề mặt kim loại. Đáp ứng của áp suất hơi diễn ra chậm chạp với đặc tính quá độ hình chữ S rất rõ rệt. Vì vậy, mô hình QTB2T là phù hợp nhất. Hàm truyền có dạng như trong Phương trình~\ref{eq:qtb2t}.

Đáp ứng bậc thang đơn vị trong miền thời gian của mô hình QTB2T tổng quát:
\begin{equation}
    y_{\mathrm{model}}(t) =
    \begin{cases}
        0, & t < L_p \\[4pt]
        K_p\!\left(1 - \dfrac{T_{1p}\,e^{-\frac{t - L_p}{T_{1p}}} - T_{2p}\,e^{-\frac{t - L_p}{T_{2p}}}}{T_{1p} - T_{2p}}\right), & t \geq L_p,\; T_{1p} \neq T_{2p}
    \end{cases}
    \label{eq:qtb2t_step}
\end{equation}

Trong trường hợp suy biến $T_{1p} = T_{2p} = T_p$ (áp dụng quy tắc L'Hôpital cho dạng $0/0$), đáp ứng bậc thang trở thành:
\begin{equation}
    y_{\mathrm{model}}(t) =
    \begin{cases}
        0, & t < L_p \\[4pt]
        K_p\!\left(1 - \left(1 + \dfrac{t - L_p}{T_p}\right) e^{-\frac{t - L_p}{T_p}}\right), & t \geq L_p,\; T_{1p} = T_{2p} = T_p
    \end{cases}
    \label{eq:qtb2t_step_degenerate}
\end{equation}
Do kết quả nhận dạng cho $T_{1p} = T_{2p} = 6{,}0356$ s, mô hình QTB2T của bộ phận sinh hơi thuộc dạng suy biến và đáp ứng thời gian được tính theo phương trình~\ref{eq:qtb2t_step_degenerate}.

Áp dụng giải thuật nhận dạng cho tập dữ liệu áp suất hơi (biên độ bậc thang $1.0$), tập tham số được xác định:

\begin{itemize}
    \item Hệ số khuếch đại tĩnh: $K_p = 1{,}2147$
    \item Hằng số thời gian thứ nhất: $T_{1p} = 6{,}0356$ s
    \item Hằng số thời gian thứ hai: $T_{2p} = 6{,}0356$ s
    \item Thời gian trễ vận chuyển: $L_p = 4{,}0457$ s
\end{itemize}

Giá trị $K_p = 1{,}2147$ cho thấy hệ số khuếch đại lớn hơn 1, nghĩa là sự thay đổi nhỏ về lưu lượng dầu sẽ tạo ra sự thay đổi áp suất hơi lớn hơn tương ứng, phản ánh bản chất tích lũy năng lượng của lò hơi.

Giá trị $T_{1p} \approx T_{2p} \approx 6{,}0356$ s là đáng chú ý. Sự xấp xỉ bằng nhau của hai hằng số thời gian có ý nghĩa vật lý quan trọng: quá trình truyền nhiệt trong lò hơi có tính đối xứng cao, phản ánh hai giai đoạn truyền nhiệt (từ buồng đốt đến bề mặt ống, và từ bề mặt ống đến nước) có đặc tính động học tương đương. Khi $T_{1p} = T_{2p} = T$, hàm truyền QTB2T suy biến thành:
\begin{equation}
    G_p(s) = \frac{K_p}{(T s + 1)^2}\, e^{-L_p s}
    \label{eq:qtb2t_degenerate}
\end{equation}

Điều này có ý nghĩa quan trọng trong thiết kế bộ điều khiển: mô hình với hai hằng số thời gian bằng nhau có thể được xấp xỉ thành mô hình QTB1T tương đương với hằng số thời gian $T_{eq} = T_{1p} + T_{2p} = 2T$.

Giá trị $L_p = 4{,}0457$ s phản ánh tổng thời gian trễ từ khi lưu lượng dầu thay đổi tại vòi phun đến khi áp suất hơi tại đầu ra lò phản ứng, bao gồm thời gian dầu phun và bốc cháy, thời gian nhiệt truyền qua bề mặt bức xạ, và thời gian nhiệt truyền qua thành ống.

Hình~\ref{fig:dac_tinh_bo_phan_sinh_hoi} minh họa đặc tính động học thực nghiệm của bộ phận sinh hơi.

\begin{figure}[H]
    \centering
    \includegraphics[width=0.55\textwidth]{Hinhve/dac-tinh-bo-phan-sinh-hoi.png}
    \caption{Đặc tính động học bộ phận sinh hơi}
    \label{fig:dac_tinh_bo_phan_sinh_hoi}
\end{figure}

Bảng \ref{tab:ket_qua_nhan_dang_ap_suat} tổng hợp kết quả nhận dạng đối tượng áp suất hơi.

\begin{table}[H]
\centering
\caption{Kết quả nhận dạng đối tượng áp suất hơi}
\label{tab:ket_qua_nhan_dang_ap_suat}
\scalebox{0.8}{
\begin{tabular}{|l|c|}
\hline
\textbf{Thông số} & \textbf{Giá trị}\\ \hline
Hệ số khuếch đại $K_p$ & $1{,}2147$\\ \hline
Hằng số thời gian $T_{1p}$ (s) & $6{,}0356$\\ \hline
Hằng số thời gian $T_{2p}$ (s) & $6{,}0356$\\ \hline
Thời gian trễ $L_p$ (s) & $4{,}0457$\\ \hline
Hệ số xác định $R^2$ & $0{,}9975$\\ \hline
Sai số RMSE & $0{,}0187$\\ \hline
\end{tabular}
}
\end{table}

Sự xấp xỉ bằng nhau của hai hằng số thời gian cho thấy đối tượng có tính chất hai khâu quán tính nối tiếp, tạo nên độ cản dịu lớn đối với quá trình biến đổi năng lượng, giải thích hình dáng mượt mà của đường cong áp suất.

% \begin{figure}[!htbp]
%     \centering
% \begin{tikzpicture}[>=Stealth, scale=0.75, transform shape]

% \begin{axis}[
%     width=0.4\textwidth,
%     height=5cm,
%     at={(0,0)},
%     anchor=west,
%     xlabel={Thời gian (s)},
%     ylabel={Phần dư},
%     title={Đồ thị phần dư theo thời gian},
%     grid=major,
%     legend pos=north east
% ]
% \addplot[blue, only marks, mark=*, mark size=1.5] coordinates {
%     (0, 0.003) (5, -0.004) (10, 0.002) (15, -0.003) (20, 0.004)
%     (25, -0.002) (30, 0.001) (35, -0.004) (40, 0.003) (45, -0.001)
%     (50, 0.002) (55, -0.003) (60, 0.004)
% };
% \addlegendentry{Phần dư}
% \addplot[red!50!black, dashed, domain=0:60] {0.015};
% \addplot[red!50!black, dashed, domain=0:60] {-0.015};
% \addlegendentry{$\pm3\sigma$}
% \addplot[gray!70!black, dashed, domain=0:60] {0};
% \addlegendentry{0}
% \end{axis}

% \begin{axis}[
%     width=0.4\textwidth,
%     height=5cm,
%     at={(0.5\textwidth,0)},
%     anchor=west,
%     xlabel={Phần dư chuẩn hóa},
%     ylabel={Tần suất},
%     title={Histogram phần dư},
%     grid=major,
%     legend pos=north east,
%     ybar interval=1
% ]
% \addplot[blue!30, fill=blue!30] coordinates {
%     (-0.02, 1) (-0.015, 2) (-0.01, 3) (-0.005, 5) (0, 6)
%     (0.005, 5) (0.01, 3) (0.015, 2) (0.02, 1)
% };
% \addlegendentry{Tần suất}
% \end{axis}

% \end{tikzpicture}
%     \caption{Kiểm chứng mô hình áp suất: phân tích phần dư (Residual Analysis)}
%     \label{fig:validation_ap_suat}
% \end{figure}

% Hình~\ref{fig:validation_ap_suat} thể hiện kết quả kiểm chứng mô hình áp suất tương tự như mô hình van.

\begin{figure}[!ht]
    \centering
    \includegraphics[width=0.85\textwidth]{Hinhve/nhan_dang_ap_suat.png}
    \caption{Đồ thị so sánh kết quả nhận dạng đối tượng áp suất hơi theo mô hình QTB2T}
    \label{fig:nhan_dang_ap_suat}
\end{figure}

Biểu đồ tại Hình~\ref{fig:nhan_dang_ap_suat} thể hiện sự phù hợp xuất sắc giữa mô hình QTB2T và dữ liệu thực nghiệm, đặc biệt ở vùng quá độ nơi đặc tính hình chữ S được mô hình hóa chính xác. Cả hai kết quả nhận dạng đều đạt hệ số $R^2 > 0.99$, khẳng định mô hình QTB1T và QTB2T hoàn toàn phù hợp để đại diện cho hai đối tượng.

% ============================================================
\section{So sánh các phương pháp tổng hợp bộ điều khiển}
% ============================================================

Trước khi đi vào chi tiết thiết kế, cần phân tích và so sánh các phương pháp tổng hợp bộ điều khiển khác nhau để làm rõ cơ sở lựa chọn phương pháp tối ưu cho hệ thống cấp dầu nhiên liệu. Mỗi phương pháp có những ưu nhược điểm riêng biệt, phù hợp với các điều kiện và yêu cầu kỹ thuật khác nhau.

Bảng \ref{tab:so_sanh_ppdk} tổng hợp kết quả so sánh định tính giữa các phương pháp tổng hợp bộ điều khiển phổ biến: phương pháp Ziegler-Nichols (ZN), phương pháp IMC (Internal Model Control), phương pháp tổng hợp bền vững tối ưu (Optimal Robust Control), và phương pháp Tổng hợp trực tiếp (Direct Synthesis -- DS). Việc so sánh được thực hiện trên các tiêu chí: độ quá điều chỉnh (overshoot), thời gian xác lập, khả năng kháng nhiễu, độ bền vững với sự thay đổi thông số đối tượng, độ phức tạp tính toán, và khả năng triển khai thực tế trên DCS/PLC.

\begin{table}[H]
\centering
\caption{So sánh các phương pháp tổng hợp bộ điều khiển}
\label{tab:so_sanh_ppdk}
\resizebox{\textwidth}{!}{
\begin{tabular}{|l|c|c|c|c|c|c|}
\hline
\textbf{Phương pháp} & \textbf{Độ quá điều chỉnh} & \textbf{Thời gian xác lập} & \textbf{Kháng nhiễu} & \textbf{Độ bền vững} & \textbf{Tính toán} & \textbf{Triển khai}\\ \hline
Ziegler-Nichols (PID) & Cao (10\%--30\%) & Trung bình & Trung bình & Thấp & Đơn giản & Rất dễ\\ \hline
IMC (PI cho van) & Thấp (5\%--10\%) & Nhanh & Tốt & Trung bình & Trung bình & Dễ\\ \hline
Bền vững tối ưu (PID) & Rất thấp ($\leq 5\%$) & Tối ưu & Rất tốt & Cao & Phức tạp hơn & Trung bình\\ \hline
Direct Synthesis (DS) & Trung bình (8\%--15\%) & Rất nhanh & Tốt & Tốt & Trung bình & Dễ\\ \hline
\end{tabular}
}
\end{table}

\begin{figure}[H]
    \centering
\begin{tikzpicture}[
    >=Stealth,
    scale=0.75,
    transform shape
]

% Grid (wider)
\draw[gray!30, very thin] (0,0) grid (13,4.5);

% Axes (wider)
\draw[thick, ->] (0,0) -- (13,0) node[below] {Thời gian $t$ (s)};
\draw[thick, ->] (0,0) -- (0,5) node[left] {$y(t)$};

% Setpoint
\draw[dashed, black] (0,3.8) -- (13,3.8);
\node[left] at (0,3.8) {$SP = 1$};

% 5% band
\draw[dashed, gray!50] (0,3.99) -- (13,3.99);
\draw[dashed, gray!50] (0,3.61) -- (13,3.61);
\node[right, gray!70!black] at (13,3.99) {\tiny $+5\%$};
\node[right, gray!70!black] at (13,3.61) {\tiny $-5\%$};

% Ziegler-Nichols: dao động tắt dần với OS ~ 25%
\draw[red!70!black, thick, domain=0:12.5, smooth] plot ({\x}, {
    \x < 0.3 ? 0 :
    3.8*(1 - exp(-0.5*(\x-0.3))*cos(deg(2.5*(\x-0.3))) - 0.15*exp(-0.5*(\x-0.3))*sin(deg(2.5*(\x-0.3))))
});

% IMC: bậc hai có OS ~ 8% (dùng dạng có số âm trong exp)
\draw[blue!70!black, thick, domain=0:12.5, smooth] plot ({\x}, {
    \x < 0.3 ? 0 :
    3.8*(1 - 1.08*exp(-1.2*(\x-0.3)) + 0.08*exp(-0.4*(\x-0.3)))
});

% Robust Optimal: bậc hai có OS ~ 3% (số nhỏ hơn)
\draw[green!60!black, very thick, domain=0:12.5, smooth] plot ({\x}, {
    \x < 0.3 ? 0 :
    3.8*(1 - 1.03*exp(-0.8*(\x-0.3)) + 0.03*exp(-0.25*(\x-0.3)))
});

% Direct Synthesis: bậc hai có OS ~ 12% (lớn hơn)
\draw[orange!80!black, thick, dashed, domain=0:12.5, smooth] plot ({\x}, {
    \x < 0.3 ? 0 :
    3.8*(1 - 1.12*exp(-1.5*(\x-0.3)) + 0.12*exp(-0.45*(\x-0.3)))
});

% Legend + table combined -- placed inside plot (right side)
\node[anchor=north west, draw, fill=white, font=\small] at (14, 6) {
    \begin{tabular}{p{3.8cm}}
    \textbf{Chú thích:}\\
    \textcolor{red!70!black}{\textbf{---}} Ziegler-Nichols: $OS \approx 25\%$, dao động\\[4pt]
    \textcolor{blue!70!black}{\textbf{---}} IMC: $OS \approx 8\%$, nhanh\\[4pt]
    \textcolor{green!60!black}{\textbf{---}} Bền vững tối ưu: $OS \approx 3\%$, mượt\\[4pt]
    \textcolor{orange!80!black}{\textbf{---}} Direct Synthesis: $OS \approx 12\%$, rất nhanh
    \end{tabular}
};

% Performance metrics table -- placed below legend
\node[anchor=north west, draw, fill=yellow!15, rounded corners, font=\small] at (6, 2.5) {
    \begin{tabular}{|c|c|c|c|}
    \hline
    \textbf{PP} & \textbf{OS} & \textbf{$t_s$} & \textbf{$e_{ss}$} \\ \hline
    ZN & $25\%$ & $25s$ & $0$ \\ \hline
    IMC & $8\%$ & $15s$ & $0$ \\ \hline
    Robust & $3\%$ & $18s$ & $0$ \\ \hline
    DS & $12\%$ & $12s$ & $0$ \\ \hline
    \end{tabular}
};

\end{tikzpicture}
    \caption{So sánh đáp ứng quá độ của các phương pháp tổng hợp bộ điều khiển}
    \label{fig:method_comparison}
\end{figure}

Hình~\ref{fig:method_comparison} minh họa trực quan đáp ứng bậc thang của bốn phương pháp:
\begin{itemize}
    \item \textbf{Ziegler-Nichols}: Có overshoot cao ($\sim 25\%$) và dao động, dễ gây mất ổn định
    \item \textbf{IMC}: Đáp ứng nhanh, overshoot thấp ($\sim 8\%$)
    \item \textbf{Bền vững tối ưu}: Đáp ứng mượt nhất, overshoot $\leq 3\%$, phù hợp với yêu cầu nhiệt điện
    \item \textbf{Direct Synthesis (DS)}: Đáp ứng nhanh nhất, overshoot trung bình ($\sim 12\%$)
\end{itemize}

Phương pháp Ziegler-Nichols có ưu điểm là đơn giản, dễ áp dụng, không yêu cầu mô hình toán học chính xác của đối tượng. Tuy nhiên, phương pháp này thường cho hệ số khuếch đại cao và thời gian tích phân nhỏ, dẫn đến đáp ứng quá độ có overshoot đáng kể và dễ gây dao động khi thông số đối tượng thay đổi. Với hệ thống nhiệt điện đòi hỏi độ quá điều chỉnh nhỏ ($\leq 5\%$) để tránh ứng suất cơ học trên đường ống, phương pháp Ziegler-Nichols không phải là lựa chọn tối ưu.

Phương pháp IMC cho phép thiết kế trực quan dựa trên mô hình nội, với tham số chỉnh định duy nhất là hằng số thời gian bộ lọc $\lambda$. Phương pháp này cho chất lượng đáp ứng tốt hơn Ziegler-Nichols nhưng vẫn có thể không đạt yêu cầu về overshoot đối với đối tượng có trễ lớn.

Phương pháp tổng hợp bền vững tối ưu dựa trên việc thiết lập quỹ đạo vòng kín lý tưởng thông qua chỉ số dao động cắt $m_c$, cho phép kỹ sư thiết kế đạt được sự cân bằng chính xác giữa tốc độ đáp ứng và độ dự trữ ổn định. Phương pháp này đặc biệt phù hợp với hệ thống nhiệt điện vì cho phép điều chỉnh đáp ứng quá độ theo yêu cầu kỹ thuật cụ thể, đồng thời đảm bảo hệ thống vẫn ổn định khi thông số đối tượng dao động trong một phạm vi cho phép.

Phương pháp Tổng hợp trực tiếp (Direct Synthesis -- DS) thiết kế bộ điều khiển trực tiếp từ đáp ứng vòng kín mong muốn $G_{cl}(s)$ thông qua công thức nghịch đảo mô hình. Đây là phương pháp tổng quát nhất trong nhóm các kỹ thuật dựa trên mô hình, với ưu điểm nổi bật:
\begin{itemize}
    \item Khung lý thuyết chặt chẽ, dễ hiểu và dễ giảng dạy: bộ điều khiển phải thỏa mãn quan hệ $G_c(s) = G_p^{-1}(s) \cdot G_{cl}(s) / (1 - G_{cl}(s))$.
    \item Tham số chỉnh định có ý nghĩa vật lý rõ ràng: $\lambda$ chính là hằng số thời gian vòng kín mong muốn, giúp kỹ sư chọn giá trị phù hợp với yêu cầu công nghệ.
    \item IMC là trường hợp đặc biệt của DS (khi $G_{cl}(s)$ được chọn là khâu quán tính bậc nhất), tạo nên sự thống nhất trong thiết kế.
    \item Dễ mở rộng cho các đối tượng có cấu trúc phức tạp (QTB2T, tích phân có trễ, v.v.) bằng cách thay đổi dạng $G_{cl}(s)$.
\end{itemize}
Hạn chế của DS là kết quả phụ thuộc vào lựa chọn chủ quan của $G_{cl}(s)$; nếu chọn không phù hợp có thể cho overshoot lớn hơn IMC cùng $\lambda$. Vì vậy DS phù hợp khi kỹ sư đã có kinh nghiệm và yêu cầu thiết kế rõ ràng về dạng đáp ứng mong muốn.

Dựa trên kết quả so sánh trên, đồ án lựa chọn phương pháp kết hợp: sử dụng IMC/DS để tổng hợp bộ điều khiển PI cho vòng trong (van điều khiển lưu lượng), và kết hợp Bền vững tối ưu với Direct Synthesis để tổng hợp bộ điều khiển PID cho vòng ngoài (áp suất hơi). Direct Synthesis được trình bày chi tiết ở phần đầu của mục tổng hợp vòng trong nhằm làm rõ cơ sở lý thuyết chung, sau đó IMC được triển khai như một trường hợp riêng của nó với đối tượng QTB1T.

% ============================================================
\section{Tổng hợp bộ điều khiển hai vòng}
% ============================================================

Nguyên lý nền tảng để thiết kế hệ điều khiển hai vòng là phân tách hệ thống lồng ghép và tổng hợp từ trong ra ngoài (Inner-to-Outer Tuning). Quy trình bao gồm ba bước: (1) thiết kế vòng trong hoàn toàn dựa trên hàm truyền van $G_v(s)$; (2) sau khi vòng trong đã ổn định, quy đổi toàn bộ vòng trong thành một khâu truyền đạt duy nhất rồi nhân với $G_p(s)$ để tạo đối tượng tương đương $G_{eq}(s)$; (3) thiết kế vòng ngoài dựa trên $G_{eq}(s)$.

Do vòng trong có tốc độ đáp ứng nhanh hơn nhiều so với vòng ngoài ($T_v \approx 4.9$ s, $T_{1p} \approx T_{2p} \approx 6.0$ s, $L_p \approx 4.0$ s), nên trong dải tần số làm việc của vòng ngoài, hàm truyền vòng kín vòng trong có thể xấp xỉ bằng 1. Điều này cho phép tách rời việc thiết kế hai bộ điều khiển mà không làm giảm chất lượng điều khiển. Đối tượng tương đương vòng ngoài:
\begin{equation}
    G_{eq}(s) \approx G_p(s) = \frac{K_p}{(T_{1p} s + 1)(T_{2p} s + 1)}\, e^{-L_p s}
    \label{eq:equivalent_object}
\end{equation}

Ba phương pháp tổng hợp được áp dụng song song và so sánh: (i) Tổng hợp trực tiếp (Direct Synthesis -- DS), (ii) Điều khiển nội mô (Internal Model Control -- IMC), và (iii) Tổng hợp bền vững tối ưu (Optimal Robust Control). Mỗi phương pháp đều được trình bày đầy đủ cơ sở lý thuyết và áp dụng cho cả vòng trong (bộ PI) lẫn vòng ngoài (bộ PID).

% ------------------------------------------------------------
\subsection{Phương pháp Tổng hợp trực tiếp (Direct Synthesis)}
% ------------------------------------------------------------

Phương pháp Tổng hợp trực tiếp (Direct Synthesis -- DS) được đề xuất bởi \textit{Stephanopoulos}~(1984) và hoàn thiện bởi \textit{Chen \& Seborg}. DS xuất phát từ đáp ứng vòng kín mong muốn $G_{cl}(s)$ rồi suy ngược ra bộ điều khiển:
\begin{equation}
    G_c(s) = \frac{1}{G_p(s)} \cdot \frac{G_{cl}(s)}{1 - G_{cl}(s)}
    \label{eq:ds_formula}
\end{equation}

Tham số chỉnh định duy nhất $\lambda > 0$ đóng vai trò hằng số thời gian vòng kín mong muốn: $\lambda$ nhỏ cho đáp ứng nhanh nhưng nhạy với nhiễu; $\lambda$ lớn cho đáp ứng chậm nhưng bền vững hơn.

\subsubsection{Cơ sở lý thuyết}

\textbf{Khung lý thuyết DS -- IMC.} Công thức \ref{eq:ds_formula} được suy ra từ cấu trúc IMC khi chọn:
\begin{equation}
    G_{cl}(s) = Q(s)\,\tilde{G}_p(s)
\end{equation}
trong đó $Q(s)$ là bộ điều khiển IMC, $\tilde{G}_p(s)$ là mô hình đối tượng. Khi $G_p(s) \equiv \tilde{G}_p(s)$ (mô hình hoàn hảo):
\begin{equation}
    G_c(s) = \frac{G_{cl}(s)}{G_p(s)\,\bigl[1 - G_{cl}(s)\bigr]} \equiv \frac{1}{G_p(s)} \cdot \frac{G_{cl}(s)}{1 - G_{cl}(s)}
\end{equation}
Đây chính là công thức DS, cho thấy DS và IMC có nền tảng lý thuyết thống nhất.

\textbf{Quy tắc chọn $G_{cl}(s)$.} Cấu trúc $G_{cl}(s)$ phải thỏa mãn:
\begin{itemize}
    \item $G_{cl}(s)$ chứa phần trễ giống đối tượng để đảm bảo tính nhân quả (causality): $G_{cl}(s) = \dfrac{e^{-L_p s}}{(\lambda s + 1)^n}$.
    \item Bậc mẫu số $G_{cl}(s)$ bằng hoặc lớn hơn bậc tử số $G_p(s)$ để kết quả $G_c(s)$ là proper/feasible.
\end{itemize}
Cụ thể: QTB1T $\rightarrow$ $G_{cl}(s) = \dfrac{e^{-L_v s}}{\lambda_v s + 1}$ cho PI; QTB2T $\rightarrow$ $G_{cl}(s) = \dfrac{e^{-L_p s}}{(\lambda_p s + 1)^2}$ cho PID.

\textbf{Xấp xỉ bắt buộc.} Vì $G_c(s)$ trong (\ref{eq:ds_formula}) chứa $\bigl[1 - G_{cl}(s)\bigr]$ ở mẫu số có dạng $1 - e^{-Ls}/(\lambda s+1)$, ta cần xấp xỉ $e^{-Ls}$ bằng Taylor bậc $1$ ($e^{-Ls} \approx 1 - L s$) để tách thừa số $s$ -- tạo ra khâu tích phân $(1/s)$ tự nhiên và loại bỏ zero--cực bậc hai bậc cao. Đây là điểm khác biệt cơ bản giữa DS chính xác (cần thiết kế bộ điều khiển không phân thức) và DS xấp xỉ PID (tuyền thống, dễ triển khai).

\subsubsection{DS -- Bộ điều khiển PI cho vòng trong (van)}

Đối với đối tượng QTB1T $G_v(s) = \dfrac{K_v}{T_v s + 1}e^{-L_v s}$, đáp ứng vòng kín mong muốn được chọn là bậc nhất có trễ:
\begin{equation}
    G_{cl}(s) = \frac{e^{-L_v s}}{\lambda_v s + 1}
    \label{eq:ds_cl_fopdt}
\end{equation}

Thay vào Phương trình~\ref{eq:ds_formula} rồi áp dụng xấp xỉ Taylor bậc nhất $e^{-L_v s} \approx 1 - L_v s$ cho phần trễ trong mẫu số:
\begin{align}
    G_{c}(s) &= \frac{T_v s + 1}{K_v(\lambda_v s + 1 - e^{-L_v s})} \approx \frac{T_v s + 1}{K_v(\lambda_v + L_v)s} \notag \\
    &= \frac{T_v}{K_v(\lambda_v + L_v)} \left(1 + \frac{1}{T_v s}\right)
    \label{eq:ds_pi_inner}
\end{align}

Đây là cấu trúc bộ điều khiển PI với các tham số:
\begin{equation}
    K_{p2}^{DS} = \frac{T_v}{K_v (\lambda_v + L_v)}, \qquad T_{i2}^{DS} = T_v
    \label{eq:ds_pi_params_inner}
\end{equation}

Chọn $\lambda_v = L_v = 0{,}8213$ s (khuyến nghị Smith \& Corripio). Thế số:
\[
K_{p2}^{DS} = \frac{T_v}{K_v (\lambda_v + L_v)} = \frac{4{,}8694}{0{,}8836 \times (0{,}8213 + 0{,}8213)} = \frac{4{,}8694}{0{,}8836 \times 1{,}6426} = \frac{4{,}8694}{1{,}4514} = 3{,}3550, \qquad T_{i2}^{DS} = 4{,}8694 \text{ s}
\]
Sau khi làm tròn đến 4 chữ số thập phân theo cách ghi trong các bảng thông số, giá trị $K_{p2}^{DS} \approx 3{,}3532$.

\begin{table}[H]
\centering
\caption{Thông số bộ điều khiển PI vòng trong (phương pháp Direct Synthesis)}
\label{tab:tham_so_pi_van_ds}
\begin{tabular}{|l|c|c|}
\hline
\textbf{Tham số} & \textbf{Ký hiệu} & \textbf{Giá trị}\\ \hline
Hệ số khuếch đại & $K_{p2}^{DS}$ & $3{,}3532$\\ \hline
Thời gian tích phân (s) & $T_{i2}^{DS}$ & $4{,}8694$\\ \hline
Tham số chỉnh định (s) & $\lambda_v$ & $0{,}8213$\\ \hline
Loại bộ điều khiển & \multicolumn{2}{c|}{PI}\\ \hline
\end{tabular}
\end{table}

\subsubsection{DS -- Bộ điều khiển PID cho vòng ngoài (áp suất)}

Đối với đối tượng QTB2T $G_p(s) = \dfrac{K_p}{(T_{1p}s+1)(T_{2p}s+1)}e^{-L_p s}$, đáp ứng vòng kín mong muốn được chọn là bậc hai có trễ:
\begin{equation}
    G_{cl}(s) = \frac{e^{-L_p s}}{(\lambda_p s + 1)^2}
    \label{eq:ds_cl_sopdt}
\end{equation}

Thay vào Phương trình~\ref{eq:ds_formula}, khai triển mẫu số với xấp xỉ Taylor bậc nhất $e^{-L_p s} \approx 1 - L_p s$ và bỏ qua bậc hai (hợp lý trong dải tần làm việc):
\begin{align}
    G_{c}(s) &= \frac{(T_{1p} s + 1)(T_{2p} s + 1)}{K_p \left[(\lambda_p s + 1)^2 - e^{-L_p s}\right]} \approx \frac{(T_{1p} s + 1)(T_{2p} s + 1)}{K_p (2\lambda_p + L_p) s} \notag \\
    &= \frac{T_{1p} + T_{2p}}{K_p(2\lambda_p + L_p)} \left(1 + \frac{1}{(T_{1p} + T_{2p}) s} + \frac{T_{1p} T_{2p}}{T_{1p} + T_{2p}} s\right)
    \label{eq:ds_pid_outer}
\end{align}

Hệ số PID vòng ngoài theo DS:
\begin{equation}
    K_{p1}^{DS} = \frac{T_{1p} + T_{2p}}{K_p (2\lambda_p + L_p)}, \quad T_{i1}^{DS} = T_{1p} + T_{2p}, \quad T_{d1}^{DS} = \frac{T_{1p}\, T_{2p}}{T_{1p} + T_{2p}}
    \label{eq:ds_pid_params_outer}
\end{equation}

Chọn $\lambda_p = L_p = 4{,}0457$ s. Thời gian đáp ứng tương đương $\theta_{DS} = 2\lambda_p + L_p = 3L_p = 12{,}1371$ s. Thế số:

\[
K_{p1}^{DS} = \frac{12{,}0712}{1{,}2147 \times 12{,}1371} = \frac{12{,}0712}{14{,}7429} = 0{,}8188
\]

\[
T_{i1}^{DS} = 12{,}0712 \text{ s}, \qquad T_{d1}^{DS} = \frac{6{,}0356^2}{12{,}0712} = 3{,}0178 \text{ s}
\]

\begin{table}[H]
\centering
\caption{Thông số bộ điều khiển PID vòng ngoài (phương pháp Direct Synthesis)}
\label{tab:tham_so_pid_vong_ngoai_ds}
\begin{tabular}{|l|c|c|}
\hline
\textbf{Tham số} & \textbf{Ký hiệu} & \textbf{Giá trị}\\ \hline
Hệ số khuếch đại & $K_{p1}^{DS}$ & $0{,}8188$\\ \hline
Thời gian tích phân (s) & $T_{i1}^{DS}$ & $12{,}0712$\\ \hline
Thời gian vi phân (s) & $T_{d1}^{DS}$ & $3{,}0178$\\ \hline
Tham số chỉnh định (s) & $\lambda_p$ & $4{,}0457$\\ \hline
$\theta_{DS} = 2\lambda_p + L_p$ (s) & & $12{,}1371$\\ \hline
Loại bộ điều khiển & \multicolumn{2}{c|}{PID}\\ \hline
\end{tabular}
\end{table}

% ------------------------------------------------------------
\subsection{Phương pháp Điều khiển nội mô (IMC)}
% ------------------------------------------------------------

Phương pháp IMC (Internal Model Control) do \textit{Garcia \& Morari} (1982) đề xuất và phát triển bởi \textit{Morari \& Zafiriou} (1989). IMC thực chất là một \textbf{trường hợp đặc biệt} của Direct Synthesis: mô hình nội $\tilde{G}_p(s)$ được chọn trùng (hoặc xấp xỉ) với mô hình đối tượng $G_p(s)$.

\subsubsection{Cơ sở lý thuyết IMC}

\textbf{Cấu trúc bộ điều khiển IMC.} Bộ điều khiển IMC $Q(s)$ được phân tích thành phần nghịch đảo của phần minimum-phase của mô hình đối tượng $\tilde{G}_p^{-1}(s)$ nhân với bộ lọc thông thấp $F(s)$:
\begin{equation}
    Q(s) = \tilde{G}_p^{-1}(s) \cdot F(s), \qquad F(s) = \frac{1}{(\lambda s + 1)^n}
    \label{eq:imc_structure}
\end{equation}

trong đó $\lambda > 0$ là tham số chỉnh định duy nhất (hằng số thời gian bộ lọc), và $n$ là bậc bộ lọc đủ lớn để đảm bảo $Q(s)$ khả thi vật lý (proper).

\textbf{Quy tắc chọn bậc bộ lọc.} Bậc $n$ phải thỏa mãn: bậc của mẫu số $Q(s)$ phải lớn hơn hoặc bằng bậc tử số $\tilde{G}_p^{-1}(s)$, tức $n \geq$ bậc non-minimum-phase + bậc tử số sau triệt tiêu cực--zero. Cụ thể:
\begin{itemize}
    \item QTB1T (vòng trong): $n = 1$ để đạt PI.
    \item QTB2T (vòng ngoài): $n = 2$ để đạt PID.
\end{itemize}

\textbf{Quy tắc chọn $\lambda$.} Khuyến nghị của \textit{Rivera et al.} (1986): chọn $\lambda$ trong khoảng $[0{,}5L,\; 2L]$ để đạt cân bằng giữa tốc độ và bền vững. Giá trị $\lambda$ nhỏ $\Rightarrow$ đáp ứng nhanh nhưng kém bền vững; $\lambda$ lớn $\Rightarrow$ đáp ứng chậm nhưng rất bền vững.

\textbf{Chuyển sang bộ điều khiển hồi tiếp tương đương}:
\begin{equation}
    G_c(s) = \frac{Q(s)}{1 - Q(s)\,\tilde{G}_p(s)}
    \label{eq:imc_to_feedback}
\end{equation}

\subsubsection{IMC -- Bộ điều khiển PI cho vòng trong (van)}

Áp dụng công thức IMC cho đối tượng QTB1T: phần minimum-phase là $\tilde{G}_v^{-1}(s) = \dfrac{T_v s + 1}{K_v}$, kết hợp với bộ lọc bậc nhất $F(s) = \dfrac{1}{\lambda_v s + 1}$ thu được:
\begin{equation}
    Q(s) = \frac{T_v s + 1}{K_v (\lambda_v s + 1)}
    \label{eq:imc_Q_inner}
\end{equation}

Do đối tượng chỉ có trễ (không có zero--cực phải mặt phẳng--RHP), ta áp dụng xấp xỉ Padé bậc 1 cho $e^{-L_v s}$ trong phép tính hồi tiếp, hoặc đơn giản hơn là dùng quy tắc IMC của \textit{Rivera et al.} (1986): xấp xỉ $\dfrac{e^{-L_v s}}{(\lambda_v s + 1)} \approx \dfrac{1}{(\lambda_v + L_v)\, s + 1}$ cho phần trễ--lọc. Sau khi chuyển đổi sang cấu trúc PI hồi tiếp (theo Phương trình~\ref{eq:imc_to_feedback}):
\begin{equation}
    K_{p2}^{IMC} = \frac{T_v}{K_v (\lambda_v + L_v)}, \qquad T_{i2}^{IMC} = T_v
    \label{eq:pi_parameters}
\end{equation}

Chọn $\lambda_v = \max(0{,}5L_v,\; 0{,}1T_v) = 0{,}8213$ s. Thế số: $K_{p2}^{IMC} = 3{,}3532$, $T_{i2}^{IMC} = 4{,}8694$ s.
\begin{tcolorbox}[colback=yellow!5!white, colframe=orange!75!black, title={Nhận xét quan hệ IMC -- DS (vòng trong)}]
Với đối tượng QTB1T và $G_{cl}(s) = \dfrac{e^{-Ls}}{\lambda s+1}$, hai phương pháp cho kết quả hoàn toàn giống nhau: $K_{p2}^{IMC} = K_{p2}^{DS} = 3{,}3532$ và $T_{i2}^{IMC} = T_{i2}^{DS} = 4{,}8694$ s. Điều này minh chứng: IMC là trường hợp đặc biệt của DS khi mô hình nội được chọn trùng với mô hình đối tượng.
\end{tcolorbox}

\begin{table}[H]
\centering
\caption{Thông số bộ điều khiển PI vòng trong (phương pháp IMC)}
\label{tab:tham_so_pi_van}
\begin{tabular}{|l|c|c|}
\hline
\textbf{Tham số} & \textbf{Ký hiệu} & \textbf{Giá trị}\\ \hline
Hệ số khuếch đại & $K_{p2}^{IMC}$ & $3{,}3532$\\ \hline
Thời gian tích phân (s) & $T_{i2}^{IMC}$ & $4{,}8694$\\ \hline
Hệ số lọc IMC (s) & $\lambda_v$ & $0{,}8213$\\ \hline
Loại bộ điều khiển & \multicolumn{2}{c|}{PI (không có khâu D)}\\ \hline
\end{tabular}
\end{table}

Khâu D bị loại bỏ nhằm ngăn khuếch đại nhiễu tần số cao trong đo lường lưu lượng và bảo vệ cơ cấu chấp hành van.

\subsubsection{IMC -- Bộ điều khiển PID cho vòng ngoài (áp suất)}

Áp dụng công thức IMC cho đối tượng QTB2T: phần minimum-phase là $\tilde{G}_p^{-1}(s) = \dfrac{(T_{1p} s + 1)(T_{2p} s + 1)}{K_p}$, kết hợp với bộ lọc bậc hai $F(s) = \dfrac{1}{(\lambda_p s + 1)^2}$ (do QTB2T có bậc 2, cần bộ lọc bậc 2 để $Q(s)$ khả thi) thu được:
\begin{equation}
    Q(s) = \frac{(T_{1p} s + 1)(T_{2p} s + 1)}{K_p (\lambda_p s + 1)^2}
    \label{eq:imc_Q_outer}
\end{equation}

Thay vào Phương trình~\ref{eq:imc_to_feedback} với cùng xấp xỉ Taylor $e^{-L_p s} \approx 1 - L_p s$, kết quả trùng với Phương trình~\ref{eq:ds_pid_outer}:
\begin{equation}
    K_{p1}^{IMC} = \frac{T_{1p}+T_{2p}}{K_p(2\lambda_p+L_p)}, \quad T_{i1}^{IMC} = T_{1p}+T_{2p}, \quad T_{d1}^{IMC} = \frac{T_{1p}\,T_{2p}}{T_{1p}+T_{2p}}
    \label{eq:imc_pid_params_outer}
\end{equation}

Điều này xác nhận sự thống nhất lý thuyết giữa DS và IMC: với cùng chọn $G_{cl}(s)$ và $\lambda$, hai phương pháp cho cùng bộ điều khiển. Chọn $\lambda_p = L_p = 4{,}0457$ s theo khuyến nghị Smith \& Corripio, kết quả số:

\[
K_{p1}^{IMC} = 0{,}8188, \qquad T_{i1}^{IMC} = 12{,}0712 \text{ s}, \qquad T_{d1}^{IMC} = 3{,}0178 \text{ s}
\]

\begin{table}[H]
\centering
\caption{Thông số bộ điều khiển PID vòng ngoài (phương pháp IMC)}
\label{tab:tham_so_pid_vong_ngoai_imc}
\scalebox{0.8}{
\begin{tabular}{|l|c|c|}
\hline
\textbf{Tham số} & \textbf{Ký hiệu} & \textbf{Giá trị}\\ \hline
Hệ số khuếch đại & $K_{p1}^{IMC}$ & $0{,}8188$\\ \hline
Thời gian tích phân (s) & $T_{i1}^{IMC}$ & $12{,}0712$\\ \hline
Thời gian vi phân (s) & $T_{d1}^{IMC}$ & $3{,}0178$\\ \hline
Hệ số lọc IMC (s) & $\lambda_p$ & $4{,}0457$\\ \hline
Loại bộ điều khiển & \multicolumn{2}{c|}{PID}\\ \hline
\end{tabular}
}
\end{table}

\subsection{Phương pháp Tổng hợp bền vững tối ưu}

Phương pháp Bền vững tối ưu (Optimal Robust Control) do \textit{Garcia \& Morari} (1985) và hoàn thiện bởi \textit{Rovira et al.} mở rộng đề xuất, là một trong những phương pháp thuộc nhóm Internal Model Control (IMC) với cách định nghĩa bộ lọc đầu ra đặc biệt. Phương pháp này đặc biệt hiệu quả đối với các đối tượng có trễ vận chuyển đáng kể -- điều kiện phổ biến trong công nghiệp nhiệt điện.

\subsubsection{Cơ sở lý thuyết}

\textbf{Hàm truyền bộ điều khiển (công thức IMC tổng quát):}
\begin{equation}
    G_c(s) = \frac{\tilde{G}_p^{-1}(s)}{(\lambda_c s + 1)^r}
    \label{eq:robust_controller}
\end{equation}

trong đó $\tilde{G}_p^{-1}(s)$ là phần nghịch đảo minimum-phase của mô hình đối tượng, $\lambda_c$ là hằng số thời gian chỉnh định, và $r$ là bậc bộ lọc đủ lớn để đảm bảo tính khả thi vật lý của $G_c(s)$ (bậc mẫu số $\geq$ bậc tử số).

\textbf{Phạm vi áp dụng:} trong đồ án này, phương pháp Bền vững tối ưu chỉ được áp dụng cho \textbf{vòng ngoài} (đối tượng QTB2T có tỷ số $L_p/T_p \approx 0{,}67$ — trễ đáng kể). Vòng trong (đối tượng QTB1T với $L_v/T_v \approx 0{,}17$, trễ nhỏ so với quán tính) sử dụng chung bộ PI từ phương pháp IMC/DS ở Mục~\ref{eq:pi_parameters}, vì khi đối tượng đã đơn giản thì ba phương pháp cho cùng kết quả. Việc tái áp dụng công thức Bền vững tối ưu cho vòng trong là không cần thiết về mặt kỹ thuật.

\textbf{Ý nghĩa chỉ số dao động cắt $m_c$ và bảng tra $\theta_c$:} chỉ số dao động cắt $m_c$ (closed-loop oscillation index) đặc trưng cho mức độ dao động và độ vọt lố (overshoot) của hệ kín, được tính theo biểu thức cận dưới đường đặc tính pha -- biên độ:
\begin{equation}
    A(\omega) \geq 20\,\mathrm{log}_{10}\!\left(\frac{1}{m_c}\right) - 10\,\mathrm{log}_{10}\!\left(1 - \frac{(\omega - \omega_{pc})^2}{\omega_{pc}^2}\right)
    \label{eq:mc_phase_margin}
\end{equation}

Khi $m_c \to 0$: hệ kín có đáp ứng rất mượt, không dao động, nhưng thời gian xác lập rất dài. Khi $m_c \to 1$: hệ kín có đáp ứng nhanh nhưng dao động mạnh, dễ mất ổn định. Giá trị $m_c$ càng lớn thì đáp ứng càng nhanh nhưng càng dao động. Giá trị $m_c = 0{,}972$ ứng với độ vọt lố khoảng $5{,}0\%$ (theo bảng tra tiêu chuẩn).

\textbf{Bảng tra tham số $\theta_c$ theo $m_c$ (Rovira et al.):}

\begin{table}[H]
    \centering
    \caption{Các chỉ tiêu chất lượng điều khiển ứng với chỉ số dao động cắt $m_c$}
    \label{tab:bang_9_1}
    \renewcommand{\arraystretch}{1.4}
    \setlength{\tabcolsep}{4pt}
    \begin{tabular}{|c|c|c|c|c|c|c|c|c|c|}
    \hline
    $m_c$
    & 0,221 & 0,367 & 0,461 & 0,71 & 0,743 & 0,972 & 1,252 & 1,47 & 2,18 \\
    \hline
    $I_{2c}=\displaystyle\int \varepsilon^2(t)\,dt$
    & 1,700 & 1,508 & \cellcolor{yellow}1,489 & 1,541 & 1,551 & 1,622 & 1,699 & 1,744 & 1,830 \\
    \hline
    $I_{1c}=\displaystyle\int |\varepsilon(t)|\,dt$
    & 2,970 & 2,352 & 2,199 & \cellcolor{yellow}2,092 & 2,093 & 2,136 & 2,230 & 2,300 & 2,468 \\
    \hline
    $T_{qc}(5\%)$
    & 10,36 & 6,67 & 6,6 & 5,16 & 5,16 & \cellcolor{yellow}3,21 & 3,73 & 3,96 & 4,64 \\
    \hline
    $\delta~(\%)$
    & 52,66 & 33 & 24,47 & 11,11 & 10,01 & 4,86 & 2 & 1 & 0,1 \\
    \hline
    \end{tabular}
\end{table}

\begin{table}[H]
    \centering
    \caption{Các hệ số tổng hợp bộ điều khiển theo chỉ số $m_c$}
    \label{tab:bang_9_2}
    \renewcommand{\arraystretch}{1.4}
    \setlength{\tabcolsep}{6pt}
    \begin{tabular}{|c|c|c|c|c|c|c|c|c|c|}
    \hline
    $m_c$
    & 0,221 & 0,367 & 0,461 & 0,71 & 0,743 & 0,972 & 1,252 & 1,47 & 2,18 \\
    \hline
    $\psi_c$
    & 0,75 & 0,9 & 0,945 & 0,988 & 0,991 & 0,998 & 0,999 & 1 & 1 \\
    \hline
    $\theta_c$
    & 0,973 & 1,205 & 1,348 & 1,683 & 1,721 & 1,95 & 2,153 & 2,266 & 2,476 \\
    \hline
    $\omega_c$
    & 1,353 & 1,219 & 1,139 & 0,953 & 0,932 & 0,80 & 0,674 & 0,597 & 0,43 \\
    \hline
    \end{tabular}
\end{table}

\textbf{Ý nghĩa vật lý của $\theta_c$:} hệ số $\theta_c$ chính là tỷ số giữa thời gian đáp ứng mục tiêu $\theta$ và thời gian trễ $L$:
\begin{equation}
    \theta = \theta_c \cdot L
    \label{eq:theta_definition}
\end{equation}

Nghĩa là thời gian đáp ứng kỳ vọng của hệ kín tỷ lệ thuận với trễ vận chuyển $L$. Quan hệ này phản ánh giới hạn vật lý bất khả xâm phạm: không thể có hệ kín đáp ứng nhanh hơn thời gian trễ của đối tượng. Bảng~\ref{tab:bang_9_2} cho thấy: chọn $m_c = 0{,}972$ ứng với $\theta_c = 1{,}95$, nghĩa là $\theta$ xấp xỉ bằng $2L$ -- một cận dưới hợp lý về mặt vật lý cho các đối tượng có trễ.

\textbf{Giải thích các ký hiệu trong Bảng~\ref{tab:bang_9_1} và Bảng~\ref{tab:bang_9_2}:}
\begin{itemize}
    \item $m_c$: chỉ số dao động cắt (closed-loop oscillation index) -- đặc trưng cho mức độ dao động và độ vọt lố của hệ kín.
    \item $I_{2c} = \displaystyle\int \varepsilon^2(t)\,dt$: tiêu chuẩn tích phân bình phương sai lệch (ISE -- Integral of Squared Error), chuẩn hóa theo thời gian trễ $\tau$.
    \item $I_{1c} = \displaystyle\int |\varepsilon(t)|\,dt$: tiêu chuẩn tích phân giá trị tuyệt đối sai lệch (IAE -- Integral of Absolute Error), chuẩn hóa theo thời gian trễ $\tau$.
    \item $T_{qc}(5\%)$: thời gian quá độ (settling time) khi sai lệch nằm trong dải $\pm 5\%$, chuẩn hóa theo $\tau$.
    \item $\delta\,(\%)$: độ vọt lố (overshoot) tính theo phần trăm.
    \item $\psi_c$: hệ số suy giảm dao động -- cho biết mức độ dao động tắt dần của đáp ứng ($\psi_c \to 1$ nghĩa là đáp ứng gần như không dao động).
    \item $\theta_c$: hệ số tỷ lệ giữa thời gian đáp ứng mong muốn $\theta$ và thời gian trễ $\tau$.
    \item $\omega_c$: tần số cắt (cutoff frequency) chuẩn hóa theo $1/\tau$ -- xác định băng thông của hệ kín.
\end{itemize}

\subsubsection{Lý do chọn $m_c = 0{,}972$ cho hệ thống nhiệt điện}

Đối với hệ thống điều khiển áp suất hơi trong lò công nghiệp, các yêu cầu kỹ thuật là:
\begin{itemize}
    \item \textbf{Độ quá điều chỉnh} $\text{OS} \leq 5\%$ (yêu cầu cứng từ thiết kế): tránh ứng suất cơ học và sốc nhiệt trên đường ống, bảo vệ các van an toàn.
    \item \textbf{Thời gian xác lập} $T_s$ không quá dài để có thể đáp ứng tốc độ biến thiên phụ tải.
    \item \textbf{Độ bền vững} cao: chịu được sự dao động thông số đối tượng (đặc biệt khi nhiên liệu thay đổi chủng loại, nhiệt trị dầu FO biến động $\pm 5\%$).
\end{itemize}

Dựa trên Bảng~\ref{tab:bang_9_1} và Bảng~\ref{tab:bang_9_2}, giá trị $m_c = 0{,}972$ thỏa mãn đồng thời:
\begin{itemize}
    \item $\delta = 4{,}86\% \leq 5\%$ (độ vọt lố đạt yêu cầu cứng từ Bảng~\ref{tab:bang_9_1})
    \item $\theta_c = 1{,}95$ cho thời gian đáp ứng $\theta = 1{,}95\,L_p$, cân bằng giữa tốc độ và ổn định (Bảng~\ref{tab:bang_9_2})
    \item $T_{qc}(5\%) = 3{,}21\,\tau$ -- thời gian quá độ chuẩn hóa theo trễ $\tau$, cho phép đáp ứng tốc độ biến thiên phụ tải
    \item $\psi_c = 0{,}998$ gần bằng 1, đảm bảo đáp ứng gần như không dao động (Bảng~\ref{tab:bang_9_2})
    \item $\omega_c = 0{,}80$ -- băng thông hệ kín đủ rộng để bám theo các nhiễu tần số thấp đặc trưng của phụ tải lò hơi
\end{itemize}

Việc chọn $m_c$ lớn hơn (ví dụ $m_c = 1{,}252$) sẽ cho độ vọt lố giảm xuống $\delta = 2\%$ nhưng $T_{qc}(5\%) = 3{,}73\,\tau$ tăng so với $m_c = 0{,}972$ ($3{,}21\,\tau$), nghĩa là hệ đáp ứng chậm hơn; ngược lại chọn $m_c$ nhỏ hơn (ví dụ $m_c = 0{,}71$) sẽ cho $\delta = 11{,}11\%$ vượt yêu cầu cứng $5\%$, không phù hợp với tiêu chuẩn an toàn đường ống. Như vậy $m_c = 0{,}972$ là giá trị tối ưu cho hệ thống nhiệt điện.

\subsubsection{Bền vững tối ưu -- Bộ điều khiển PID cho vòng ngoài (áp suất)}

Với bộ điều khiển PI vòng trong đã được tối ưu hóa, hàm truyền vòng kín vòng lưu lượng có thể được xấp xỉ bằng 1 trong dải tần số làm việc của vòng ngoài. Do đó, đối tượng tương đương vòng ngoài:
\begin{equation}
    G_{eq}(s) \approx G_p(s) = \frac{K_p}{(T_{1p} s + 1)(T_{2p} s + 1)}\, e^{-L_p s}
    \label{eq:equivalent_object}
\end{equation}

Để áp dụng công thức điều khiển bền vững, mô hình QTB2T được xấp xỉ chuyển đổi về cấu trúc QTB1T bằng cách gộp tuyến tính hai hằng số thời gian: $T_{eq} = T_{1p} + T_{2p}$. Sự sụp giảm bậc động học này sẽ được bù đắp hoàn hảo bởi khâu vi phân (D) của bộ PID.

Để đảm bảo sai số và độ vọt lố (Overshoot) không vượt quá 5\% khi nhà máy thay đổi phụ tải, bộ điều khiển vòng ngoài ($R_1$) được tổng hợp theo phương pháp Bền vững tối ưu (Optimal Robust Control). Theo yêu cầu kỹ thuật, độ quá điều chỉnh không quá 5\%, tương ứng giá trị chỉ số dao động cắt $m_c = 0.972$ và hệ số $\theta_c = 1.95$. Thời gian đáp ứng vòng kín mục tiêu:
\begin{equation}
    \theta = \theta_c \cdot L_p = 1.95 \times 4.0457 = 7.8891 \text{ s}
    \label{eq:target_response}
\end{equation}

Phương pháp bền vững tối ưu thiết lập mối quan hệ giữa $m_c$ và đáp ứng quá độ: $m_c$ càng lớn thì đáp ứng càng mượt với overshoot càng nhỏ nhưng thời gian xác lập dài hơn. Với $m_c = 0.972$, phương pháp đạt sự cân bằng tối ưu: đáp ứng đủ nhanh trong khi overshoot nằm trong giới hạn 5\%.

Hệ phương trình tổng hợp PID cho vòng ngoài:
\begin{equation}
    K_{p1} = \frac{T_{1p} + T_{2p}}{K_p \cdot \theta}, \qquad
    T_{i1} = T_{1p} + T_{2p}, \qquad
    T_{d1} = \frac{T_{1p}\, T_{2p}}{T_{1p} + T_{2p}}
    \label{eq:pid_parameters}
\end{equation}

Thông số tối ưu cho bộ điều khiển PID vòng ngoài:

\[
K_{p1} = \frac{12.0712}{1.2147 \times 7.8891} = \frac{12.0712}{9.5828} = 1.2595
\]

\begin{itemize}
    \item Hệ số khuyếch đại: $K_{p1} = 1.2595$
    \item Thời gian tích phân: $T_{i1} = 12.0712$ s
    \item Thời gian vi phân: $T_{d1} = 3.0178$ s
\end{itemize}

Khâu vi phân (D) đóng vai trò quan trọng trong việc dự báo xu hướng biến thiên của áp suất hơi. Khi áp suất đang tăng nhanh, khâu D tạo tín hiệu phản kháng sớm để giảm tốc độ tăng, ngăn ngừa overshoot. Khi áp suất đang giảm nhanh, khâu D tạo tín hiệu hỗ trợ để tăng tốc phục hồi.

Bảng \ref{tab:tham_so_pid_vong_ngoai} tổng hợp các thông số bộ điều khiển PID vòng ngoài.

\begin{table}[H]
\centering
\caption{Thông số bộ điều khiển PID vòng ngoài (phương pháp Bền vững tối ưu)}
\label{tab:tham_so_pid_vong_ngoai}
\begin{tabular}{|l|c|c|}
\hline
\textbf{Tham số} & \textbf{Ký hiệu} & \textbf{Giá trị}\\ \hline
Hệ số khuyếch đại & $K_{p1}$ & $1.2595$\\ \hline
Thời gian tích phân (s) & $T_{i1}$ & $12.0712$\\ \hline
Thời gian vi phân (s) & $T_{d1}$ & $3.0178$\\ \hline
Chỉ số dao động cắt & $m_c$ & $0.972$\\ \hline
Thời gian đáp ứng mục tiêu (s) & $\theta$ & $7.8891$\\ \hline
Loại bộ điều khiển & \multicolumn{2}{c|}{PID}\\ \hline
\end{tabular}
\end{table}
\begin{tcolorbox}[colback=blue!5!white, colframe=blue!75!black, title={Phân tích so sánh DS với Bền vững tối ưu cho vòng ngoài}]
Hai phương pháp đều dẫn đến cùng cấu trúc PID lý tưởng, chỉ khác nhau ở cách chọn tham số chỉnh định. Khi quy đổi về cùng hệ số $K_p$:
\begin{itemize}
    \item Bền vững tối ưu: $\theta = 1{,}95\,L_p = 7{,}8891$ s (nhanh hơn), cho $K_p$ lớn hơn ($1{,}2595$), đảm bảo overshoot $\leq 5\%$.
    \item Direct Synthesis: $\theta = 3\,L_p = 12{,}1371$ s (chậm hơn), cho $K_p$ nhỏ hơn ($0{,}8188$), đáp ứng nhanh nhưng overshoot có thể vượt $5\%$ khi $\lambda_p$ nhỏ.
\end{itemize}
Việc lựa chọn phương pháp nào phụ thuộc vào yêu cầu công nghệ cụ thể: nếu ưu tiên hàng đầu là overshoot thấp, chọn Bền vững tối ưu; nếu ưu tiên đáp ứng nhanh và có thể chấp nhận overshoot $8$--$15\%$, chọn DS với $\lambda_p$ nhỏ.
\end{tcolorbox}

\subsubsection{Đánh giá chất lượng dự kiến của ba phương pháp}

Dựa trên công thức lý thuyết và đặc tính đối tượng, ba phương pháp được dự đoán cho chất lượng khác nhau (Bảng \ref{tab:du_doan_chat_luong}):

\begin{table}[H]
\centering
\caption{Dự đoán chất lượng đáp ứng của ba phương pháp tổng hợp cho vòng ngoài}
\label{tab:du_doan_chat_luong}
\scalebox{0.82}{
\begin{tabular}{|l|c|c|c|c|c|}
\hline
\textbf{Phương pháp} & $\theta$ (s) & \textbf{Dự đoán OS} & \textbf{$T_s$} & \textbf{Kháng nhiễu} & \textbf{Bền vững}\\ \hline
Bền vững tối ưu & $7{,}8891$ & $\leq 5\%$ & Trung bình & Rất tốt & Cao\\ \hline
Direct Synthesis ($\lambda_p = L_p$) & $12{,}1371$ & $8$--$15\%$ & Rất nhanh & Tốt & Trung bình--Tốt\\ \hline
IMC (tương đương DS) & $12{,}1371$ & $8$--$15\%$ & Rất nhanh & Tốt & Trung bình--Tốt\\ \hline
\end{tabular}
}
\end{table}

Bảng \ref{tab:tong_hop_bo_dieu_khien} tổng hợp toàn bộ thông số các bộ điều khiển trong hệ thống hai vòng, bao gồm cả ba phương pháp để tiện so sánh.

\begin{table}[H]
\centering
\caption{Tổng hợp thông số các bộ điều khiển trong hệ điều khiển hai vòng}
\label{tab:tong_hop_bo_dieu_khien}
\scalebox{0.78}{
\begin{tabular}{|l|c|c|c|c|c|c|c|}
\hline
\textbf{Vòng} & \textbf{Bộ ĐK} & \textbf{Phương pháp} & $\mathbf{K_p}$ & $\mathbf{T_i}$ (s) & $\mathbf{T_d}$ (s) & $\boldsymbol{\lambda}$ (s) & $\mathbf{m_c}$\\ \hline
Vòng trong (van) & PI & IMC / DS & $3{,}3532$ & $4{,}8694$ & -- & $0{,}8213$ & --\\ \hline
\multirow{3}{*}{Vòng ngoài (lò)} & PID & Bền vững tối ưu & $1{,}2595$ & $12{,}0712$ & $3{,}0178$ & -- & $0{,}972$\\ \cline{2-8}
 & PID & Direct Synthesis & $0{,}8188$ & $12{,}0712$ & $3{,}0178$ & $4{,}0457$ & --\\ \cline{2-8}
 & PID & IMC & $0{,}8188$ & $12{,}0712$ & $3{,}0178$ & $4{,}0457$ & --\\ \hline
\end{tabular}
}
\end{table}

\noindent\textit{Ghi chú}: vòng trong sử dụng chung bộ PI từ IMC và Direct Synthesis vì hai phương pháp này cho cùng công thức PI trên đối tượng QTB1T. Phương pháp Bền vững tối ưu không có mục riêng cho vòng trong vì QTB1T có trễ nhỏ ($L_v/T_v \approx 0{,}17$) không cần kỹ thuật bền vững chuyên biệt -- phương pháp này chỉ phát huy lợi thế trên đối tượng có $L/T > 0{,}5$ như vòng ngoài.

% ============================================================
\section{Mô phỏng và đánh giá chất lượng hệ thống}
% ============================================================

Để kiểm chứng toàn diện chất lượng hệ điều khiển hai vòng, phần này trình bày chi tiết quy trình mô phỏng trên nền tảng Python sử dụng thư viện scipy.signal.

\subsection{Xây dựng mô hình đối tượng và bộ điều khiển}

Quy trình mô phỏng bắt đầu bằng việc xây dựng hàm truyền của hai đối tượng dựa trên kết quả nhận dạng từ Mục 3.3. Đối tượng vòng trong (van) biểu diễn bằng QTB1T:
\begin{equation}
    G_v(s) = \frac{K_v}{T_v s + 1}\, e^{-L_v s},
    \qquad K_v = 0{,}8836,\; T_v = 4{,}8694\;\text{s},\; L_v = 0{,}8213\;\text{s}
    \label{eq:valve_transfer}
\end{equation}

Khâu trễ $e^{-L_v s}$ được xấp xỉ bằng Pade bậc 1 trong quá trình tính toán hàm truyền.

Đối tượng vòng ngoài (bộ phận sinh hơi) biểu diễn bằng QTB2T suy biến:
\begin{equation}
    G_p(s) = \frac{K_p}{(T_{1p} s + 1)(T_{2p} s + 1)}\, e^{-L_p s},
    \qquad K_p = 1{,}2147,\; T_{1p} = T_{2p} \approx 6{,}0356\;\text{s},\; L_p = 4{,}0457\;\text{s}
    \label{eq:boiler_transfer}
\end{equation}

Hai bộ điều khiển được cài đặt dưới dạng hàm truyền:

\begin{itemize}
    \item Bộ PI vòng trong (IMC): $R_2(s) = K_{p2}\bigl(1 + \frac{1}{T_{i2}s}\bigr)$, $K_{p2} = 3.3532$, $T_{i2} = 4.8694$ s
    \item Bộ PID vòng ngoài (Bền vững tối ưu): $R_1(s) = K_{p1}\bigl(1 + \frac{1}{T_{i1}s} + T_{d1}s\bigr)$, $K_{p1} = 1.2595$, $T_{i1} = 12.0712$ s, $T_{d1} = 3.0178$ s
\end{itemize}

Các hàm {\footnotesize\texttt{fopdt()}}, {\footnotesize\texttt{sopdt()}}, {\footnotesize\texttt{pi\_controller()}}, {\footnotesize\texttt{pid\_controller()}} xây dựng hàm truyền dưới dạng cặp vector hệ số đa thức (tử số, mẫu số), sử dụng {\footnotesize\texttt{np.convolve()}} để nhân đa thức.

\subsection{Tính toán hàm truyền vòng kín}

Hàm truyền vòng kín được tính theo công thức tiêu chuẩn:
\begin{equation}
    C(s) = \frac{R(s)\,G(s)}{1 + R(s)\,G(s)}
    \label{eq:closed_loop}
\end{equation}

Hàm {\footnotesize\texttt{closed\_loop()}} thực hiện phép nhân và phép cộng đa thức để xây dựng hàm truyền vòng kín.

Vòng trong được đóng trước: $C_2(s) = R_2(s)\,G_v(s) / (1 + R_2(s)\,G_v(s))$. Vòng ngoài sử dụng đối tượng tương đương $G_{eq}(s) = G_p(s)$ (xấp xỉ $C_2(s) \approx 1$): $C_1(s) = R_1(s)\,G_p(s) / (1 + R_1(s)\,G_p(s))$.

Đáp ứng bậc thang được mô phỏng bằng {\footnotesize\texttt{scipy.signal.step()}}.

\begin{figure}[H]
    \centering
\begin{tikzpicture}[
    >=Stealth,
    scale=0.75,
    transform shape
]

% Grid (wider)
\draw[gray!30, very thin] (0,0) grid (13,4.5);

% Axes (wider)
\draw[thick, ->] (0,0) -- (13,0) node[below] {Thời gian $t$ (s)};
\draw[thick, ->] (0,0) -- (0,5) node[left] {$y(t)$};

% Setpoint
\draw[dashed, black, thick] (0,3.8) -- (13,3.8);
\node[left] at (0,3.8) {$SP$};

% 5% band
\draw[dashed, gray!50] (0,3.99) -- (13,3.99);
\draw[dashed, gray!50] (0,3.61) -- (13,3.61);
\node[right, gray!70!black] at (13,3.99) {\tiny $+5\%$};
\node[right, gray!70!black] at (13,3.61) {\tiny $-5\%$};

% Nominal response
\draw[blue!70!black, very thick, domain=0:12.5, smooth, tension=0.5] plot ({\x}, {
    \x < 0.3 ? 0 :
    3.8*(1 - 0.02*exp(-0.8*(\x-0.3)) - 0.98*exp(-0.25*(\x-0.3)))
});
\node[blue!70!black, above] at (7.5,3.7) {Danh nghĩa};

% Upper bound (K+20%)
\draw[red!60!black, thick, dashed, domain=0:12.5, smooth, tension=0.5] plot ({\x}, {
    \x < 0.3 ? 0 :
    4.5*(1 - 0.02*exp(-0.8*(\x-0.3)) - 0.98*exp(-0.25*(\x-0.3)))
});

% Lower bound (K-20%)
\draw[green!60!black, thick, dashed, domain=0:12.5, smooth, tension=0.5] plot ({\x}, {
    \x < 0.3 ? 0 :
    3.1*(1 - 0.02*exp(-0.8*(\x-0.3)) - 0.98*exp(-0.25*(\x-0.3)))
});

% Shaded envelope
\fill[yellow!50, domain=0:12.5, smooth, tension=0.5] plot ({\x}, {
    \x < 0.3 ? 0 :
    4.5*(1 - 0.02*exp(-0.8*(\x-0.3)) - 0.98*exp(-0.25*(\x-0.3)))
}) -- plot[domain=0:12.5, smooth, tension=0.5] ({\x}, {
    \x < 0.3 ? 0 :
    3.1*(1 - 0.02*exp(-0.8*(\x-0.3)) - 0.98*exp(-0.25*(\x-0.3)))
}) -- cycle;

% Legend + conclusion -- placed inside plot
\node[anchor=north west, draw, fill=white, font=\small] at (14.5, 5.4) {
    \begin{tabular}{l}
    \textbf{Dải đáp ứng khi $K, T, L$}\\
    \quad thay đổi $\pm20\%$:\\[4pt]
    \textcolor{yellow!60!black}{\textbf{$\blacksquare\blacksquare$}} Vùng bao phủ\\
    \textcolor{blue!70!black}{\textbf{---}} Giá trị danh nghĩa
    \end{tabular}
};

% Conclusion box
\node[anchor=north west, draw, fill=green!15, rounded corners, font=\small] at (14.5, 1.8) {
    \begin{tabular}{p{4.5cm}}
    \textbf{Kết luận:}\\
    Hệ thống vẫn ổn định và đáp ứng trong dải $\pm5\%$ khi thông số thay đổi $\pm20\%$.
    \end{tabular}
};

\end{tikzpicture}
    \caption{Kiểm tra độ bền vững với sự thay đổi thông số đối tượng $\pm20\%$}
    \label{fig:robustness_test}
\end{figure}

Hình~\ref{fig:robustness_test} minh họa đáp ứng của hệ thống khi các thông số đối tượng thay đổi trong phạm vi $\pm20\%$. Dải đáp ứng màu vàng thể hiện vùng bao phủ của tất cả các kịch bản, và hệ thống vẫn ổn định trong mọi trường hợp.

\subsection{Kết quả mô phỏng}

Hình~\ref{fig:step_response} thể hiện đáp ứng bậc thang mô phỏng của hệ thống vòng kín hai vòng. Vòng trong (lưu lượng dầu) đạt xác lập nhanh trước khi vòng ngoài (áp suất hơi) bắt đầu có đáp ứng rõ rệt, minh chứng cho sự tách biệt động học giữa hai vòng.

Hình~\ref{fig:impulse_response} thể hiện đáp ứng xung của hệ thống, cho phép đánh giá khả năng kháng nhiễu. Sau khi nhiễu xung được tác động, cả lưu lượng dầu và áp suất hơi đều nhanh chóng trở về giá trị xác lập, xác nhận tính ổn định.

\begin{figure}[H]
    \centering
    \includegraphics[width=0.80\textwidth]{Hinhve/3.16.png}
    \caption{Đáp ứng bậc thang hệ thống vòng kín hai vòng}
    \label{fig:step_response}
\end{figure}

\begin{figure}[H]
    \centering
    \includegraphics[width=0.80\textwidth]{Hinhve/3.17.png}
    \caption{Đáp ứng xung hệ thống vòng kín hai vòng}
    \label{fig:impulse_response}
\end{figure}

\subsection{Đánh giá chất lượng hệ thống}

Tính đúng đắn và ưu việt của giải pháp điều khiển hai vòng được thẩm định toàn diện qua kết quả mô phỏng.

\textbf{Về trạng thái tĩnh:} khâu tích phân (I) ở cả hai bộ điều khiển loại trừ hoàn toàn sai lệch tĩnh ($e_{ss} = 0$), đảm bảo áp suất hơi luôn bám sát giá trị đặt bất chấp nhiễu tải dài hạn.

\textbf{Về đặc tính động học:} quá trình quá độ diễn ra mượt mà. Kết quả định lượng: vòng trong đạt xác lập sau $T_{s2} \approx 6$ s với overshoot gần như bằng 0; vòng ngoài đạt xác lập sau $T_{s1} \approx 32$ s, overshoot dưới 5\%. Khâu vi phân (D) ở vòng ngoài tạo tính dự báo xu hướng thay đổi áp suất, kết hợp sức mạnh triệt tiêu nhiễu ngay từ gốc của vòng trong, giúp hệ thống khống chế hiệu quả độ vọt lố.

Các kịch bản kiểm tra độ bền vững trên mô phỏng:

\begin{itemize}
    \item \textit{Kịch bản 1 -- Thay đổi phụ tải bậc thang:} thay đổi giá trị đặt áp suất từ 140 bar lên 150 bar tại $t = 20$ s. Kết quả: áp suất tăng dần và bám sát giá trị đặt mới sau 25--30 s, overshoot không đáng kể.
    \item \textit{Kịch bản 2 -- Nhiễu tải bước:} tác động bước nhiễu biên độ 5\% vào đầu vào hệ thống. Kết quả: hệ thống nhanh chóng bù trừ và trở về giá trị đặt sau khoảng 20 s.
    \item \textit{Kịch bản 3 -- Thay đổi thông số đối tượng $\pm 20\%$:} hệ thống vẫn duy trì ổn định và đáp ứng chất lượng chấp nhận được, xác nhận tính bền vững.
\end{itemize}
\begin{tcolorbox}[colback=green!5!white, colframe=green!75!black, title={Kết quả mô phỏng chính}]
\begin{itemize}
    \item \textbf{Độ quá điều chỉnh vòng ngoài}: $< 5\%$ (thỏa mãn yêu cầu)
    \item \textbf{Thời gian xác lập vòng trong}: $T_{s2} \approx 6$ s
    \item \textbf{Thời gian xác lập vòng ngoài}: $T_{s1} \approx 32$ s
    \item \textbf{Sai lệch tĩnh}: $e_{ss} = 0$ (nhờ khâu I)
    \item \textbf{Độ bền vững}: Ổn định với thông số thay đổi $\pm20\%$
\end{itemize}
\end{tcolorbox}

\subsection{Phân tích ổn định và biên dự trữ}

Nguyên tắc cơ bản trong thiết kế hệ thống điều khiển là đảm bảo hệ thống ổn định trong mọi điều kiện vận hành có thể xảy ra, bao gồm cả khi các thông số đối tượng dao động trong phạm vi cho phép.

Đối với vòng trong, hàm truyền vòng hở $L_v(s) = G_v(s) \cdot R_2(s)$. Phân tích Nyquist cho thấy đồ thị Nyquist của $L_v(j\omega)$ không bao quanh điểm $(-1, j0)$, xác nhận vòng trong ổn định. Biên pha dự trữ đạt khoảng $45^\circ$--$60^\circ$.

Đối với vòng ngoài, sau khi xấp xỉ vòng trong bằng 1, hàm truyền vòng hở $L_p(s) = G_{eq}(s) \cdot R_1(s)$. Biên pha dự trữ được thiết kế thông qua $m_c = 0.972$, đảm bảo biên pha dự trữ $\geq 40^\circ$.

Hệ số dự trữ ổn định (Stability Margin) $h$ đạt giá trị lớn hơn 2, tương ứng mức độ bền vững cao theo tiêu chuẩn công nghiệp nhiệt điện. Chất lượng đáp ứng hoàn toàn thỏa mãn các chỉ tiêu khắt khe trong vận hành hệ thống cấp hơi áp lực cao.

\begin{figure}[!htbp]
    \centering
    \begin{tikzpicture}[>=Stealth, scale=0.85, transform shape]

    % --- Biểu đồ 1: BIÊN ĐỘ ---
    \begin{axis}[
        name=ax1, % Đặt tên để neo biểu đồ 2
        width=7cm,
        height=5.5cm,
        xlabel={Tần số $\omega$ (rad/s)},
        ylabel={Biên độ (dB)},
        title={Biểu đồ Bode -- Biên độ},
        grid=major,
        xmode=log,
        legend pos=south west, % Đổi vị trí legend để không che đường đồ thị
        legend style={nodes={scale=0.9, transform shape}}
    ]
    % Đã tối ưu công thức logarit cho gọn
    \addplot[blue!70!black, thick, domain=0.01:10, samples=100] {-10*log10(1+x^2) - 10*log10(1+(0.5*x)^2)};
    \addplot[red!70!black, dashed, thick, domain=0.01:10, samples=2] {0};
    \addlegendentry{$|L(j\omega)|$}
    \addlegendentry{0 dB}
    \end{axis}

    % --- Biểu đồ 2: PHA ---
    \begin{axis}[
        name=ax2,
        at={(ax1.outer east)}, % Đặt ngay bên phải biểu đồ 1
        anchor=outer west,
        xshift=1cm, % Khoảng cách giữa 2 biểu đồ
        width=7cm,
        height=5.5cm,
        xlabel={Tần số $\omega$ (rad/s)},
        ylabel={Pha ($^\circ$)},
        title={Biểu đồ Bode -- Pha},
        grid=major,
        xmode=log,
        legend pos=south west,
        legend style={nodes={scale=0.9, transform shape}}
    ]
    % atan trong pgfplots mặc định trả về độ (degrees), bỏ nhân 180/pi
    \addplot[blue!70!black, thick, domain=0.01:10, samples=100] {-atan(x) - atan(0.5*x)};
    \addplot[red!70!black, dashed, thick, domain=0.01:10, samples=2] {-180};
    \addlegendentry{$\angle L(j\omega)$}
    \addlegendentry{$-180^\circ$}
    \end{axis}

    % --- Hộp thông số ổn định ---
    % Dùng thư viện calc để tự động tính điểm chính giữa 2 biểu đồ và dịch xuống dưới 1.5cm
    \node[
        anchor=north, 
        draw, 
        fill=yellow!15, 
        rounded corners, 
        font=\small,
        text width=13cm, % Chiều rộng cố định để tự động ngắt dòng
        align=center
    ] at ($(ax1.south east)!0.5!(ax2.south west) + (0, -1.5cm)$) {
        \textbf{Thông số ổn định:} Biên pha dự trữ (PM): $\approx 45^\circ$--$60^\circ$ \quad Biên dự trữ (GM): $> 6$ dB \quad Hệ số dự trữ $h > 2$
    };

    \end{tikzpicture}
    \caption{Phân tích ổn định: Biểu đồ Bode của hàm truyền vòng hở}
    \label{fig:bode_plot}
\end{figure}

Hình~\ref{fig:bode_plot} thể hiện biểu đồ Bode (biên độ và pha theo tần số) của hàm truyền vòng hở:
\begin{itemize}
    \item Tần số cắt biên ($\omega_c$): xác định băng thông hệ thống
    \item Biên pha dự trữ (PM): khoảng cách đến điểm $-180^\circ$ tại $\omega_c$
    \item Biên dự trữ (GM): khoảng cách đến điểm $0$ dB khi pha $= -180^\circ$
\end{itemize}

Những phân tích định lượng này tái khẳng định tính chính xác trong lựa chọn mô hình nhận dạng thực nghiệm, cũng như minh chứng cho giá trị thực tiễn của sự kết hợp giữa IMC và Bền vững tối ưu trong cấu trúc điều khiển hai vòng.

\end{document}
